% !TEX TS-program = LuaLaTeX-dev
% !TEX root = ../TRENCH_DIFFEQ.tex

\chapter{First Order Equations}\label{ch:first_order}


\noindent
\textsc{In this chapter} we study
first order equations for which there are general methods of
solution.


\medskip\noindent
\zcref[S,CAP,font=\scshape]{sec:linar_first_order} deals with linear equations,
the simplest kind
of first order equations. In this section we introduce the method of
variation of parameters. The idea underlying this method will be a
unifying theme for our approach to solving many different kinds of
differential equations throughout the book.

\medskip\noindent
\zcref[S,CAP,font=\scshape]{sec:separable} deals with separable equations, the simplest
nonlinear equations. In this section we introduce the idea of implicit
and constant solutions of differential equations, and we point out
some differences between the properties of linear and nonlinear
equations.

\medskip\noindent
\zcref[S,CAP,font=\scshape]{sec:exist_unique} discusses existence and
uniqueness of solutions of nonlinear equations. Although it may seem
logical to place this section before \zcref*{sec:separable}, we
presented
\zcref*{sec:separable} first so we could have illustrative examples in
\zcref*{sec:exist_unique}.

\medskip\noindent
\zcref[S,CAP,font=\scshape]{sec:transform_separable} deals with nonlinear equations that
are not separable, but can be transformed into separable equations by a
procedure similar to variation of parameters.


\medskip\noindent
\zcref[S,CAP,font=\scshape]{sec:exact} covers exact differential equations, which are
given this name because the method for solving them uses the idea of
an exact differential from calculus.


\medskip\noindent
\zcref[S,CAP,font=\scshape]{sec:integrating_factors} deals with
equations that are not exact, but can made exact by multiplying them
by a function known called \emph{\color{blue}integrating factor}.

\newpage


\section{Linear First Order Equations}\label{sec:linar_first_order}

A first order differential equation is said to be \emph{\color{blue}linear} if it
can be written as

\begin{equation}\label{eq:2.1.1}
y'+p(x)y=f(x).
\end{equation}
A first order differential equation that can't be written like this
is \emph{\color{blue}nonlinear}. We say that \eqref{eq:2.1.1} is \emph{\color{blue}homogeneous}
if $f\equiv0$;     otherwise it's \emph{\color{blue}nonhomogeneous}. Since
$y\equiv0$ is obviously a solution of the homgeneous equation
\[
y'+p(x)y=0,
\]
we call it the \emph{\color{blue}trivial solution}. Any  other solution
is \emph{\color{blue}nontrivial}.


\begin{example}\label{example:2.1.1}
 The first order equations
\begin{align*}
x^2y'+3y&=x^2,\\
xy'-8x^2y&=\sin x,\\
xy'+(\ln x)y&=0,\\
y'&=x^2y - 2,
\end{align*}
are not in the form \eqref{eq:2.1.1}, but they are linear, since they can
be rewritten as
\begin{align*}
y'+\frac{3}{x^2}y&=1,\\
y'-8xy&=\frac{\sin x}{x},\\
y'+\frac{\ln x}{x}y&=0,\\
y'-x^2y&=-2.
\end{align*}
\end{example}

\begin{example}\label{example:2.1.2}
Here are some
nonlinear first order equations:
\begin{align*}
xy'+3y^2 & = 2x&&\text{(because $y$ is squared)},\\
yy'& =3&&\text{(because of the product $yy'$)},\\
y'+xe^y &=12&&\text{(because of $e^y$)}.
\end{align*}
\end{example}

\boxit{General Solution of a Linear First Order Equation}

To motivate a definition that we'll need, consider the
simple linear first order equation
\begin{equation}\label{eq:2.1.2}
y'=\frac{1}{x^2}.
\end{equation}
From calculus we know that $y$ satisfies this equation
if and only if
\begin{equation}\label{eq:2.1.3}
y=-\frac{1}{x}+c,
\end{equation}
where $c$ is an arbitrary constant.
We call $c$ a
\emph{\color{blue}parameter} and say that \eqref{eq:2.1.3} defines a
\emph{\color{blue}one-parameter family} of functions.
For each real number $c$, the
function defined by \eqref{eq:2.1.3} is a solution of \eqref{eq:2.1.2}
on $(-\infty,0)$ and $(0,\infty)$;     moreover, every solution of
\eqref{eq:2.1.2} on either of these intervals is of the form
\eqref{eq:2.1.3} for some choice of $c$. We say that \eqref{eq:2.1.3}
is \emph{\color{blue}the general solution} of~\eqref{eq:2.1.2}.

We'll  see that a similar situation occurs in connection with any
first order linear equation
\begin{equation}\label{eq:2.1.4}
y'+p(x)y=f(x);
\end{equation}
that is, if $p$ and $f$ are continuous on some open interval $(a,b)$ then
 there's a unique formula $y=y(x,c)$ analogous to \eqref{eq:2.1.3}
that involves $x$ and a parameter $c$ and has  the these
properties:
\begin{itemize}
\item % ()
For each fixed value of $c$, the resulting function of $x$ is a
solution of \eqref{eq:2.1.4} on $(a,b)$.
\item % ()
If $y$ is a solution of \eqref{eq:2.1.4} on $(a,b)$, then $y$ can be
obtained from the formula by choosing $c$ appropriately.
\end{itemize}
We'll call $y=y(x,c)$ the \emph{\color{blue}general solution} of
\eqref{eq:2.1.4}.

When this has been established, it will follow that
an equation of the form
\begin{equation}\label{eq:2.1.5}
P_0(x)y'+P_1(x)y=F(x)
\end{equation}
has a general solution on  any open interval $(a,b)$ on which $P_0$,
$P_1$, and $F$ are all continuous and $P_0$ has no zeros, since in
this case we can rewrite \eqref{eq:2.1.5} in the form \eqref{eq:2.1.4}
with $p=P_1/P_0$ and $f=F/P_0$, which are both continuous on $(a,b)$.

To avoid awkward wording in examples and exercises, we won't specify
the interval $(a,b)$ when we ask for the general solution of a
specific linear first order equation. Let's agree that this always
means that we want the general solution on every open interval on
which $p$ and $f$ are continuous if the equation is of the form
\eqref{eq:2.1.4}, or on which $P_0$, $P_1$, and $F$ are continuous and
$P_0$ has no zeros, if the equation is of the form \eqref{eq:2.1.5}. We
leave it to you to identify these intervals in specific examples and
exercises.


For completeness, we point out that if $P_0$, $P_1$, and $F$ are all
continuous on an open interval $(a,b)$, but $P_0$ \emph{\color{blue}does} have a
zero in $(a,b)$, then \eqref{eq:2.1.5} may fail to have a general
solution on $(a,b)$ in the sense just defined. Since this isn't  a
major point that needs to be developed in depth, we won't discuss it
further; however, see \zcref{exer:2.1.44} for an example.

\boxit{Homogeneous Linear First Order Equations}

We begin with the problem of finding the general solution of
a homogeneous linear first order equation.
The next example recalls a familiar result from calculus.


\begin{example}\label{example:2.1.3}
   Let $a$ be a constant.
\begin{alist}
\item %(a)
Find the general solution of
\begin{equation}
y'-ay=0.\label{eq:2.1.6}
\end{equation}

\item %(b)
Solve the initial value problem
\[
y'-ay=0,\quad y(x_0)=y_0.
\]
\end{alist}
\end{example}

\begin{alistSolution}
\item
You already know from calculus that if $c$ is any constant, then
$y=ce^{ax}$ satisfies \eqref{eq:2.1.6}. However, let's pretend  you've
forgotten this, and use this problem to illustrate a general method
for solving a homogeneous linear first order equation.

We know that \eqref{eq:2.1.6} has the trivial solution $y\equiv0$. Now
suppose   $y$ is a nontrivial solution of \eqref{eq:2.1.6}. Then, since
a differentiable function must be continuous, there must be some open
interval $I$ on which $y$ has no zeros. We rewrite \eqref{eq:2.1.6} as
\[
\frac{y'}{y}=a
\]
for $x$ in $I$. Integrating this shows that
\[
\ln|y|=ax+k,\mbox{\quad so \quad} |y|=e^ke^{ax},
\]
where $k$ is an arbitrary constant. Since $e^{ax}$ can never equal
zero,  $y$ has no zeros, so $y$ is either always
positive or always negative. Therefore we can rewrite $y$ as
\begin{equation}\label{eq:2.1.7}
y=ce^{ax}
\end{equation}
where
\[
c=\begin{cases}
\phantom{-}e^k&\mbox{if } y>0,\\
-e^k&\mbox{if } y<0.
\end{cases}
\]
This shows that every nontrivial solution of \eqref{eq:2.1.6} is of the
form $y=ce^{ax}$ for some nonzero constant $c$. Since setting $c=0$
yields the trivial solution,  \emph{\color{blue}all} solutions of
\eqref{eq:2.1.6} are of the form \eqref{eq:2.1.7}. Conversely,
\eqref{eq:2.1.7} is a solution of \eqref{eq:2.1.6} for every choice of $c$,
since differentiating \eqref{eq:2.1.7} yields $y'=ace^{ax}=ay$.

\item   Imposing  the initial condition $y(x_0)=y_0$
yields $y_0=ce^{ax_0}$, so $c=y_0e^{-ax_0}$ and
\[
y=y_0e^{-ax_0}e^{ax}=y_0e^{a(x-x_0)}.
\]
\zcref[S]{figure:2.1.1} show the
graphs of this function with $x_{0}=0$, $y_{0}=1$, and various values of
$a$.
\end{alistSolution}



\begin{figure}
  \centering
  \includegraphics[alt={Six solution curves corresponding to the values -4, -2.5, -1, 1, 1.5, and 2 for a.  The more negative values of a decrease more quickly.  The more positive values of a increase more quickly.},
  width=5.4in,height=3.3in,keepaspectratio]{fig020101}
\color{blue}
  \caption{Solutions of $y'-ay=0$, $y(0)=1$}
  \label{figure:2.1.1}
\end{figure}


\begin{example} \label{example:2.1.4}
\begin{alist}
\item % (a)
Find the general solution of
\begin{equation}
 xy'+y=0.\label{eq:2.1.8}
\end{equation}

\item % (b)
Solve the initial value problem
\begin{equation}
 xy'+y=0,\quad y(1)=3.\label{eq:2.1.9}
\end{equation}
\end{alist}
\end{example}


\begin{alistSolution}
\item
We rewrite \eqref{eq:2.1.8} as
\begin{equation}\label{eq:2.1.10}
y'+\frac{1}{x}y=0,
\end{equation}
where $x$ is restricted to either $(-\infty,0)$ or $(0,\infty)$. If
$y$ is a nontrivial solution of \eqref{eq:2.1.10}, there must be some
open interval I on which $y$ has no zeros. We can rewrite
\eqref{eq:2.1.10} as
\[
\frac{y'}{y}=-\frac{1}{x}
\]
for $x$ in $I$. Integrating shows that
\[
\ln|y|=-\ln|x|+k,\quad\text{so}\quad|y|=\frac{e^k}{|x|}.
\]
Since a function that satisfies the last equation can't change sign on
either $(-\infty,0)$ or $(0,\infty)$, we can rewrite this result more
simply as
\begin{equation}\label{eq:2.1.11}
y=\frac{c}{x}
\end{equation}
where
\[
c=
\begin{cases}
\phantom{-}e^k&\text{if } y>0,\\
-e^k&\text{if } y<0.
\end{cases}
\]
We've now shown that every solution of \eqref{eq:2.1.10} is given by
\eqref{eq:2.1.11} for some choice of $c$. (Even though we assumed that $y$
was nontrivial to derive \eqref{eq:2.1.11}, we can get the trivial
solution by setting $c=0$ in \eqref{eq:2.1.11}.) Conversely, any function
of the form \eqref{eq:2.1.11} is a solution of \eqref{eq:2.1.10}, since
differentiating \eqref{eq:2.1.11} yields
\[
y'=-\frac{c}{x^2},
\]
and substituting this and \eqref{eq:2.1.11} into  \eqref{eq:2.1.10} yields
\begin{align*}
y'+\frac{1}{x}y&=-\frac{c}{x^2}+\frac{1}{x}\frac{c}{x}\\
&=-\frac{c}{x^2}+\frac{c}{x^2}=0.
\end{align*}

\zcref[S]{figure:2.1.2} shows the graphs of some solutions corresponding
to various values of $c$


\item Imposing the initial condition $y(1)=3$ in
\eqref{eq:2.1.11} yields $c=3$. Therefore the solution of \eqref{eq:2.1.9} is
\[
y=\frac{3}{x}.
\]
The interval of validity of this solution is  $(0,\infty)$.
\end{alistSolution}

   \begin{figure}
  \centering
  \includegraphics[alt={Several hyperbolas y=c/x.  When c>0 the hyperbolas are in the first and third quadrants.  When c<0 the hyperbolas are in the second and fourth quadrants.},
  width=5.4in,height=3.3in,keepaspectratio]{fig020102}
\color{blue}
  \caption{Solutions of $xy'+y=0$ on $(0,\infty)$ and $(-\infty,0)$}
  \label{figure:2.1.2}
\end{figure}

The results in Examples~\ref{example:2.1.3}(a) and \ref{example:2.1.4}(b)
are special cases of the next theorem.

\begin{theorem}\label{thmtype:2.1.1}
If $p$ is continuous on $(a,b),$ then the general solution
of the homogeneous equation
\begin{equation}\label{eq:2.1.12}
y'+p(x)y=0
\end{equation}
on $(a,b)$ is
\[
y=ce^{-P(x)},
\]
where
\begin{equation}\label{eq:2.1.13}
P(x)=\int p(x)\,dx
\end{equation}
is any antiderivative of $p$ on $(a,b);$ that is$,$
\begin{equation}\label{eq:2.1.14}
P'(x)=p(x),\quad a<x<b.
\end{equation}
\end{theorem}


\begin{proof}
If $y=ce^{-P(x)}$,  differentiating $y$ and using \eqref{eq:2.1.14}
shows that
\[
y'=-P'(x)ce^{-P(x)}=-p(x)ce^{-P(x)}=-p(x)y,
\]
so $y'+p(x)y=0$; that is, $y$ is a solution of \eqref{eq:2.1.12}, for
any choice of $c$.

Now we'll show that any solution of \eqref{eq:2.1.12} can be written as
$y=ce^{-P(x)}$ for some constant $c$. The trivial solution can be
written this way, with $c=0$. Now suppose   $y$ is a nontrivial
solution. Then there's an open subinterval $I$ of $(a,b)$ on which
$y$ has no zeros. We can rewrite \eqref{eq:2.1.12} as
\begin{equation}\label{eq:2.1.15}
\frac{y'}{y}=-p(x)
\end{equation}
for $x$ in $I$. Integrating \eqref{eq:2.1.15} and recalling \eqref{eq:2.1.13}
yields
\[
\ln|y|=-P(x) + k,
\]
where $k$ is a constant. This implies that
\[
|y|=e^ke^{-P(x)}.
\]
Since $P$ is defined for all $x$ in $(a,b)$  and an exponential can
never equal zero,  we can take $I=(a,b)$, so $y$
has no zeros on %$(a,b)$
$(a,b)$, so we can rewrite the last equation as $y=ce^{-P(x)}$, where
\[
c=\begin{cases}
\phantom{-}e^k&\mbox{if } y>0\mbox{ on }(a,b),\\
-e^k&\mbox{if } y<0\mbox{ on }(a,b).
\end{cases}\qedhere
\]
\end{proof}

{
\color{blue}
\remark{Rewriting a first order differential equation so that one
side depends only on $y$ and $y'$ and the other depends only on $x$
is called \emph{\color{blue}separation of variables}. We did this in
\zcref{example:2.1.3,example:2.1.4}, and in rewriting
\eqref{eq:2.1.12} as \eqref{eq:2.1.15}. We'll apply this method to nonlinear
equations in \zcref{sec:separable}.}
}

\boxit{Linear Nonhomogeneous First Order Equations}

We'll now solve the nonhomogeneous equation
\begin{equation}\label{eq:2.1.16}
y'+p(x)y=f(x).
\end{equation}
When considering this equation  we call
\[
y'+p(x)y=0
\]
the \emph{\color{blue}complementary equation}.

We'll find solutions of \eqref{eq:2.1.16} in the form $y=uy_1$, where
$y_1$ is a nontrivial solution of the complementary equation and $u$
is to be determined. This method of using a solution of the
complementary equation to obtain solutions of a nonhomogeneous
equation is a special case of a method called \emph{\color{blue}variation of
parameters}, which you'll encounter several times in this book.
(Obviously, $u$ can't be constant, since if it were, the left side
of \eqref{eq:2.1.16} would be zero. Recognizing this, the early users of
this method viewed $u$ as a ``parameter'' that varies; hence, the name
``variation of parameters.'')

If
\[
y=uy_1,\quad\text{then}\quad y'=u'y_1+uy_1'.
\]
Substituting these expressions for $y$ and $y'$ into \eqref{eq:2.1.16}
yields
\[
u'y_1+u(y_1'+p(x)y_1)=f(x),
\]
which reduces to
\begin{equation}\label{eq:2.1.17}
u'y_1=f(x),
\end{equation}
since $y_1$ is a solution of the complementary equation; that is,
\[
y_1'+p(x)y_1=0.
\]
In the proof of \zcref{thmtype:2.2.1} we saw that $y_1$
has no zeros on an interval where $p$ is continuous. Therefore
we can divide \eqref{eq:2.1.17} through by $y_1$ to obtain
\[
u'=f(x)/y_1(x).
\]
We can integrate this (introducing a constant of integration), and
multiply  the result by $y_1$ to get the general solution of
\eqref{eq:2.1.16}.
Before turning to the formal proof of this claim, let's consider
some examples.

\begin{example}\label{example:2.1.5}
Find the general solution of
\begin{equation}\label{eq:2.1.18}
y'+2y=x^3e^{-2x}.
\end{equation}
\end{example}


By applying (a) of \zcref{example:2.1.3} with $a=-2$, we see
that $y_1=e^{-2x}$ is a solution of the  complementary equation
$y'+2y=0$. Therefore we seek solutions of \eqref{eq:2.1.18} in the form
 $y=ue^{-2x}$,  so that
\begin{equation} \label{eq:2.1.19}
 y'=u'e^{-2x}-2ue^{-2x}\quad\text{and}\quad
y'+2y=u'e^{-2x}-2ue^{-2x}+2ue^{-2x}=u'e^{-2x}.
\end{equation}
Therefore $y$ is a solution of \eqref{eq:2.1.18} if and only if
\[
u'e^{-2x}=x^3e^{-2x}\quad\text{or, equivalently,}\quad u'=x^3.
\]
Therefore
\[
u=\frac{x^4}{4}+c,
\]
and
\[
y=ue^{-2x}=e^{-2x}\left(\frac{x^4}{4}+c\right)
\]
is the general  solution of \eqref{eq:2.1.18}.

\zcref[S]{figure:2.1.3} shows a direction field and some integral curves
for \eqref{eq:2.1.18}.

\begin{figure}
  \centering
  \includegraphics[alt={A field for -1<x<1 and -2<y<2.  For x near -1, the arrows point down and to the right.  As x increases, the arrows turn to point to the right and toward y=0.  Several solution curves follow the arrows.},
  width=5.67in,height=3.3in,keepaspectratio]{fig020103}
\color{blue}
  \caption{A direction field and integral curves for $y'+2y=x^{2}e^{-2x}$}
  \label{figure:2.1.3}
\end{figure}


\begin{example}\label{example:2.1.6}
\begin{alist}
\item %(a)
Find the general solution
\begin{equation}\label{eq:2.1.20}
y'+(\cot x)y=x\csc x.
\end{equation}
\item %(b)
Solve the initial value problem
\begin{equation}\label{eq:2.1.21}
y'+(\cot x)y=x\csc x,\quad y(\pi/2)=1.
\end{equation}
\end{alist}
\end{example}

\begin{alistSolution}
\item Here $p(x)=\cot x$ and $f(x)= x\csc x$ are both
continuous except at the points $x=r\pi$, where $r$ is an integer.
Therefore we seek solutions of \eqref{eq:2.1.20} on the intervals
$\left(r\pi, (r+1)\pi \right)$.
We need a nontrival solution $y_1$ of
the complementary equation;   thus, $y_1$ must satisfy
 $y_1'+(\cot x)y_1=0$, which we rewrite as
\begin{equation}\label{eq:2.1.22}
\frac{y_1'}{y_1}=-\cot x=-\frac{\cos x}{\sin x}.
\end{equation}
Integrating this yields
\[
\ln|y_1|=-\ln|\sin x|,
\]
where we take the constant of integration to be zero since we need
only \emph{\color{blue}one} function that satisfies \eqref{eq:2.1.22}. Clearly
$y_1=1/\sin x$ is a suitable choice. Therefore we seek solutions of
\eqref{eq:2.1.20} in the form
\[
y=\frac{u}{\sin x},
\]
so that
\begin{equation} \label{eq:2.1.23}
 y'=\frac{u'}{\sin x}-\frac{u\cos x}{\sin^2x}
\end{equation}
and
\begin{align}
y'+(\cot x)y
&=\dfrac{u'}{\sin x}-\dfrac{u\cos x}{\sin^2x}+\dfrac{u\cot x}{\sin x}\notag\\
&=\dfrac{u'}{\sin x}-\dfrac{u\cos x}{\sin^2x}+\dfrac{u\cos x}{\sin^2 x}
\label{eq:2.1.24}\\
&=\dfrac{u'}{\sin x}.\notag
\end{align}
Therefore $y$ is a solution of \eqref{eq:2.1.20} if and only if
\[
u'/\sin x=x\csc x=x/\sin x\mbox{\quad or, equivalently, \quad} u'=x.
\]
Integrating this yields
\begin{equation}\label{eq:2.1.25}
u=\frac{x^2}{2}+c, \text{\quad and \quad }
y=\frac{u}{\sin x}=\frac{x^2}{2\sin x}+\frac{c}{\sin x}.
\end{equation}
is the general solution of \eqref{eq:2.1.20} on every interval
$\left(r\pi,(r+1)\pi\right)$ ($r=$integer).

\item
Imposing the initial condition $y(\pi/2)=1$ in \eqref{eq:2.1.25} yields
\[
1=\frac{\pi^2}{8}+c\mbox{\quad or\quad} c=1-\frac{\pi^2}{8}.
\]
 Thus,
\[
y=\frac{x^2}{2\sin x}+\frac{(1-\pi^2/8)}{\sin x}
\]
is a solution of \eqref{eq:2.1.21}.
The interval of validity of this solution is $(0,\pi)$;
\zcref{figure:2.1.4} shows its graph.
\end{alistSolution}

\begin{figure}
\centering
\includegraphics[alt={An increasing blue curve beginning near negative infinity as x approaches 0, crossing the x axis near (0.7,0), and finishing near positive infinity as x approaches π.},
width=5.4in,height=3in,keepaspectratio]{fig020104}
\color{blue}
\caption{Solution of
$y'+(\cot x)y=x\csc x, y(\pi/2)=1$}
  \label{figure:2.1.4}
\end{figure}

{
\color{blue}
\remark{It wasn't necessary to do the computations
\eqref{eq:2.1.23} and \eqref{eq:2.1.24} in
\zcref{example:2.1.6}, since we showed in the discussion preceding
\zcref{example:2.1.5} that if $y=uy_1$ where $y_1'+p(x)y_1=0$,
then
$y'+p(x)y=u'y_1$. We did these computations so you would see this
happen in this specific example. We recommend that you include these
``unnecessary'' computations in doing exercises, until you're confident
that you really understand the method. After that, omit them.}
}

\medskip


We summarize the method  of variation of parameters for
solving
\begin{equation} \label{eq:2.1.26}
y'+p(x)y=f(x)
\end{equation}
as follows:

\medskip
\begin{alist}[before={\color{blue}}]
\item % (a)
Find a function $y_1$  such that
\[
\frac{y_1'}{y_1}=-p(x).
\]
For convenience, take the constant of integration to be zero.
\item % (b)
Write
\begin{equation} \label{eq:2.1.27}
y=uy_1
\end{equation}
to remind yourself of what you're doing.
\item % (c)
Write $u'y_1=f$ and solve for $u'$;   thus, $u'=f/y_1$.
\item % (d
Integrate $u'$ to obtain $u$, with an arbitrary constant of
integration.
\item % (e)
Substitute $u$ into \eqref{eq:2.1.27}  to obtain $y$.
\end{alist}

\medskip
To solve an equation written as
\[
P_0(x)y'+P_1(x)y=F(x),
\]
we recommend that you divide through by $P_0(x)$ to obtain an equation
of the form  \eqref{eq:2.1.26}
and then follow this   procedure.

\boxit{Solutions in Integral Form}

Sometimes the integrals that arise in solving a linear first order
equation can't be evaluated in terms of elementary functions. In this
case the solution must be left in terms of an integral.

\begin{example}\label{example:2.1.7}
\begin{alist}
\item % (a)
Find the general solution of
\[
y'-2xy=1.
\]
\item % (b)
Solve the initial value problem
\begin{equation}\label{eq:2.1.28}
y'-2xy=1,\quad y(0)=y_0.
\end{equation}
\end{alist}
\end{example}

\begin{alistSolution}
\item
To apply variation of parameters, we need a nontrivial solution $y_1$
of  the complementary equation;   thus,
$y_1'-2xy_1=0$, which we rewrite as
\[
\frac{y_1'}{y_1}=2x.
\]
Integrating this and taking the constant of integration to be
zero yields
\[
\ln|y_1|=x^2,\mbox{\quad so \quad}|y_1|=e^{x^2}.
\]
We choose $y_1=e^{x^2}$ and seek solutions of
 \eqref{eq:2.1.28} in the form $y=ue^{x^2}$, where
\[
u'e^{x^2}=1,\mbox{\quad so \quad} u'=e^{-x^2}.
\]
Therefore
\[
u=c+\int e^{-x^2}dx,
\]
 but we can't simplify the integral on the right because
there's  no elementary function with  derivative equal to $e^{-x^2}$.
Therefore the best available form for the  general solution of
\eqref{eq:2.1.28} is
\begin{equation}\label{eq:2.1.29} y=ue^{x^2}= e^{x^2}\left(c+\int
e^{-x^2}dx\right). \end{equation}

\item
  Since the initial condition in \eqref{eq:2.1.28}
is imposed at $x_0=0$, it is convenient to rewrite \eqref{eq:2.1.29} as
\[
y=e^{x^2}\left(c+\int^x_0 e^{-t^2} dt\right),\mbox{\quad since
\quad} \int_0^0e^{-t^2}\,dt=0.
\]
 Setting $x=0$ and $y=y_0$ here shows that
$c=y_0$.  Therefore the solution of the initial value problem is
\begin{equation}\label{eq:2.1.30}
y=e^{x^2}\left(y_0 +\int^x_0 e^{-t^2}dt\right).
\end{equation}
For a given value of $y_0$ and each fixed $x$, the integral on the
right can be evaluated by numerical methods. An alternate procedure is
to apply the numerical integration procedures discussed in
Chapter~3 directly to the initial value problem
\eqref{eq:2.1.28}.
\zcref[S]{figure:2.1.5} shows graphs of
 of \eqref{eq:2.1.30} for several values of $y_0$.
\end{alistSolution}


\begin{figure}
  \centering
  \includegraphics[alt={Ten solution curves that begin at x=0 and go to the right, turning upward as x increases.},
  width=5.4in,height=3.3in,keepaspectratio]{fig020105}
\color{blue}
  \caption{Solutions of $y'-2xy=1$, $y(0)=y_{0}$}
  \label{figure:2.1.5}
\end{figure}

\boxit{An Existence and Uniqueness Theorem}

The method of variation of parameters leads to this theorem.

\begin{theorem}\label{thmtype:2.1.2}
Suppose $p$ and $f$ are continuous on an open interval $(a,b),$  and let $y_1$
be any nontrivial solution of the complementary equation
\[
y'+p(x)y=0
\]
on $(a,b)$. Then:
\begin{alist}
\item % (a)
The general solution of the nonhomogeneous equation
\begin{equation}\label{eq:2.1.31}
y'+p(x)y=f(x)
\end{equation}
 on  $(a,b)$ is
\begin{equation}\label{eq:2.1.32}
y=y_1(x)\left(c +\int f(x)/y_1(x)\,dx\right).
\end{equation}
\item % (b)
If $x_0$ is an arbitrary point in $(a,b)$ and $y_0$ is an arbitrary
real number$,$ then the initial value problem
\[
y'+p(x)y=f(x),\quad y(x_0)=y_0
\]
 has the unique solution
\[
y=y_1(x)\left(\frac{y_0}{y_1(x_0)} +\int^x_{x_0} \frac{f(t)}{y_1(t)}\, dt\right)
\]
on $(a,b).$
\end{alist}
\end{theorem}

\begin{proof}
\begin{alist}
\item To show that \eqref{eq:2.1.32} is the general solution of
\eqref{eq:2.1.31} on $(a,b)$, we must prove that:

\begin{rmlist}
\item % (i)
If $c$ is any constant,   the function $y$ in
\eqref{eq:2.1.32} is a solution of \eqref{eq:2.1.31} on  $(a,b)$.
\item % (ii)
  If $y$ is a solution of \eqref{eq:2.1.31} on  $(a,b)$ then $y$
is of the form  \eqref{eq:2.1.32} for some constant~$c$.
\end{rmlist}

To prove (i), we first observe that any function of the form
\eqref{eq:2.1.32} is defined on $(a,b)$, since $p$ and $f$ are continuous
on $(a,b)$. Differentiating \eqref{eq:2.1.32} yields
\[
y'=y_1'(x)\left(c +\int f(x)/y_1(x)\,
dx\right)+f(x).
\]
Since $y_1'=-p(x)y_1$, this and \eqref{eq:2.1.32} imply that
\begin{align*}
y'&=-p(x)y_1(x)\left(c +\int f(x)/y_1(x)\, dx\right)+f(x)\\
&=-p(x)y(x)+f(x),
\end{align*}
which implies that $y$ is a solution of \eqref{eq:2.1.31}.

To prove (ii), suppose $y$ is a solution of \eqref{eq:2.1.31}
on $(a,b)$. From the proof of \zcref{thmtype:2.1.1}, we know that
$y_1$ has no zeros on $(a,b)$, so the function $u=y/y_1$ is defined
on $(a,b)$. Moreover, since
\[
y'=-py+f\mbox{\quad and\quad}y_1'=-py_1,
\]
we have
\begin{align*}
u'&=\frac{y_1y'-y_1'y}{y_1^2}\\
&=\frac{y_1(-py+f)-(-py_1)y}{y_1^2}=\frac{f}{y_1}.
\end{align*}
Integrating $u'=f/y_1$ yields
\[
u=\left(c +\int f(x)/y_1(x)\, dx\right),
\]
which implies \eqref{eq:2.1.32}, since $y=uy_1$.

\item We've proved (a), where $\int f(x)/y_1(x)\,dx$ in
\eqref{eq:2.1.32}
is an arbitrary antiderivative of $f/y_1$. Now it's convenient to
choose the antiderivative that equals zero when $x=x_0$, and write the
general solution of \eqref{eq:2.1.31} as
\[
y=y_1(x)\left(c +\int^x_{x_0} \frac{f(t)}{y_1(t)}\, dt\right).
\]
Since
\[
y(x_0)= y_1(x_0)\left(c +\int^{x_0}_{x_0} \frac{f(t)}{y_1(t)}\, dt\right)=cy_1(x_0),
\]
we see that $y(x_0)=y_0$ if and only if $c=y_0/y_1(x_0)$.\qedhere
\end{alist}
\end{proof}

\exercises

\emph{In \zcref[range]{exer:2.1.1,exer:2.1.5} find the general solution.}

\begin{exerciselist}

\item(1)\label{exer:2.1.1}$y'+ay=0$ ($a$=constant)
\item(1)\label{exer:2.1.2}$y'+3x^2y=0$
\item(1)\label{exer:2.1.3}$xy'+(\ln x)y=0$
\item(1)\label{exer:2.1.4}$xy'+3y=0$

\item\label{exer:2.1.5}$x^2y'+y=0$

\exercisetext{In \zcref[range]{exer:2.1.6,exer:2.1.11} solve
the initial value problem.}

\item\label{exer:2.1.6}$y'+\left(\dfrac{1+x}{x}\right)y=0,\quad y(1)=1$

\item\label{exer:2.1.7}$xy'+\left(1+\dfrac{1}{\ln x}\right)y=0,\quad y(e)=1$

\item\label{exer:2.1.8}$xy'+(1+ x\cot x)y=0,\quad y\left(\dfrac{\pi}{2}
\right)=2$

\item\label{exer:2.1.9}$y'-\left(\dfrac{2x}{1+x^2}\right)y=0,\quad y(0)=2$

\item\label{exer:2.1.10}$y'+\dfrac{k}{x}y=0,\quad y(1)=3$, where $k$ is
constant

\item\label{exer:2.1.11}$y'+(\tan kx)y=0,\quad y(0)=2$, where $k$ is
constant

\exercisetext{In \zcref[range]{exer:2.1.12,exer:2.1.15} find
the general solution. Also, plot a direction field and some integral
curves on the rectangular region $\{-2\le x\le2,\ -2\le y\le2$\}.}

\item(1)\label{exer:2.1.12}\CGex $y'+3y=1$
\item(1)\label{exer:2.1.13}\CGex $y'+\left(\dfrac{1}{x}-1\right)y=-\dfrac{2}{x}$
\item(1)\label{exer:2.1.14}\CGex $y'+2xy=xe^{-x^2}$
\item(1)\label{exer:2.1.15}\CGex $y'+\dfrac{2x}{1+x^2}y=\dfrac{e^{-x}}{1+x^2}$

\exercisetext{In \zcref[range]{exer:2.1.16,exer:2.1.24} find
the general solution.}


\item(1)\label{exer:2.1.16}$y'+\dfrac{1}{x}y=\dfrac{7}{x^2}+3$
\item(1)\label{exer:2.1.17}$y'+\dfrac{4}{x-1}y
= \dfrac{1}{(x-1)^5}+\dfrac{\sin x}{(x-1)^4}$
\item(1)\label{exer:2.1.18}$xy'+(1+2x^2)y=x^3e^{-x^2}$
\item(1)\label{exer:2.1.19}$xy'+2y=\dfrac{2}{x^2}+1$
\item(1)\label{exer:2.1.20}$y'+(\tan x)y=\cos x$
\item(1)\label{exer:2.1.21}$(1+x)y'+2y=\dfrac{\sin x}{ 1 + x}$
\item(1)\label{exer:2.1.22}$(x-2)(x-1)y'-(4x-3)y=(x-2)^3$
\item(1)\label{exer:2.1.23}$y'+(2\sin x\cos x) y=e^{-\sin^2x}$
\item(1)\label{exer:2.1.24}$x^2y'+3xy=e^x$

\exercisetext{In \zcref[range]{exer:2.1.25,exer:2.1.29} solve the
initial value problem and sketch the graph of the solution.}

\item\label{exer:2.1.25}\CGex $y'+7y=e^{3x},\quad y(0)=0$

\item\label{exer:2.1.26}\CGex $(1+x^2)y'+4xy=\dfrac{2}{1+x^2},\quad y(0)=1$

\item\label{exer:2.1.27}\CGex $xy'+3y=\dfrac{2}{x(1+x^2)},\quad y(-1)=0$


\item\label{exer:2.1.28}\CGex
$y'+ (\cot x)y=\cos x,\quad y\left(\dfrac{\pi}{2}\right)=1$

\item\label{exer:2.1.29}\CGex $y'+\dfrac{1}{x}y=\dfrac{2}{x^2}+1,\quad y(-1)=0$


\exercisetext{In \zcref[range]{exer:2.1.30,exer:2.1.37} solve
the initial value problem.}

\item\label{exer:2.1.30}$(x-1)y'+3y=\dfrac{1}{(x-1)^3} +
\dfrac{\sin x}{(x-1)^2},\quad y(0)=1$

\item\label{exer:2.1.31}$xy'+2y=8x^2,\quad y(1)=3$

\item\label{exer:2.1.32}$xy'-2y=-x^2,\quad y(1)=1$

\item\label{exer:2.1.33}$y'+2xy=x,\quad y(0)=3$

\item\label{exer:2.1.34}$(x-1)y'+3y=\dfrac{1+(x-1)\sec^2x}{(x-1)^3},\quad y(0)=-1$

\item\label{exer:2.1.35}$(x+2)y'+4y=\dfrac{1+2x^2}{x(x+2)^3},\quad y(-1)=2$

\item\label{exer:2.1.36}$(x^2-1)y'-2xy=x(x^2-1),\quad y(0)=4$

\item\label{exer:2.1.37}$(x^2-5)y'-2xy=-2x(x^2-5),\quad y(2)=7$

\exercisetext{In \zcref[range]{exer:2.1.38,exer:2.1.42} solve
the initial value problem and leave the answer in a form involving a
definite integral. $($You can solve these problems numerically by
methods discussed in Chapter~3.$)$}

\item\label{exer:2.1.38}$y'+2xy=x^2,\quad y(0)=3$

\item\label{exer:2.1.39}$y'+\dfrac{1}{x}y=\dfrac{\sin x}{x^2},\quad y(1)=2$

\item\label{exer:2.1.40}$y'+y=\dfrac{e^{-x}\tan x}{x},\quad y(1)=0$

\item\label{exer:2.1.41}$y'+\dfrac{2x}{1+x^2}y=\dfrac{e^x}{(1+x^2)^2}, \quad y(0)=1$

\item\label{exer:2.1.42}$xy'+(x+1)y=e^{x^2},\quad y(1)=2$

\item\label{exer:2.1.43}Experiments indicate that glucose is absorbed by the body at a rate
proportional to the amount of glucose present in the bloodstream. Let
$\lambda$ denote the (positive) constant of proportionality. Now
suppose   glucose is injected into a patient's bloodstream at a
constant rate of $r$ units per unit of time. Let $G=G(t)$ be the
number of units in the patient's bloodstream at time $t>0$. Then
\[
G'=-\lambda G+r,
\]
where the first term on the right is due to the absorption of the
glucose by the patient's body and the second term is due to the
injection. Determine $G$ for $t>0$, given that $G(0)=G_0$. Also,
find $\lim_{t\to\infty}G(t)$.


\item\label{exer:2.1.44}\leavevmode\vspace{-\baselineskip}
\begin{alist}
\item % (a)
\Lex Plot a direction field and some integral curves for
\[
xy'-2y=-1\tag{A}
\]
on the rectangular region  $\{-1\le x\le 1, -.5\le y\le 1.5\}$.
What do all the integral curves have in common?
\item %(b)
Show that the general solution of (A)
on $(-\infty,0)$ and $(0,\infty)$ is
\[
y=\frac{1}{2}+cx^2.
\]
\item % (c)
Show that $y$ is a solution of (A) on
$(-\infty,\infty)$ if and only if
\[
y=\begin{cases}
\dfrac{1}{2}+c_1x^2, &x \ge 0,\\
\dfrac{1}{2}+c_2x^2, &x < 0,
\end{cases}
\]
 where $c_1$ and $c_2$ are arbitrary constants.
\item % (c)
Conclude from \part{c} that all solutions of (A) on
$(-\infty,\infty)$ are solutions of the initial value problem
\[
xy'-2y=-1,\quad y(0)=\frac{1}{2}.
\]
\item %(d)
Use \part{b} to show that if  $x_0\ne0$ and $y_0$  is arbitrary, then
  the  initial value problem
\[
xy'-2y=-1,\quad y(x_0)=y_0
\]
has infinitely many solutions on ($-\infty,\infty$). Explain why this
doesn't contradict \zcref{thmtype:2.1.1}\part{b}.
\end{alist}


\item\label{exer:2.1.45}Suppose $f$ is continuous on an open interval $(a,b)$
and $\alpha$ is a constant.
\begin{alist}
\item %(a)
Derive a formula for the solution of the initial value
problem
\renewcommand{\theequation}{\Alph{equation}}
\[
y'+\alpha y=f(x),\quad y(x_0)=y_0,\tag{A}
\]
where $x_0$ is in $(a,b)$ and $y_0$ is an arbitrary real number.

\item %(b)
Suppose $(a,b)=(a,\infty)$, $\alpha > 0$ and
$\displaystyle{\lim_{x\to\infty} f(x)=L}$. Show that if $y$ is the
solution of (A), then $\displaystyle{\lim_{x\to
\infty} y(x)=L/\alpha}$.
\end{alist}

\item\label{exer:2.1.46}Assume that all functions in this exercise are defined on
a common interval $(a,b)$.
\begin{alist}
\item %(a)
Prove:  If $y_1$ and $y_2$ are solutions of
\[
y'+p(x)y=f_1(x)
\]
and
\[
y'+p(x)y=f_2(x)
\]
respectively, and $c_1$ and $c_2$ are constants, then
$y=c_1y_1+c_2y_2$ is a solution of
\[
y'+p(x)y=c_1f_1(x)+c_2f_2(x).
\]
(This is the \emph{principle of superposition}.)

\item %(b)
Use \part{a} to show that if $y_1$ and $y_2$ are solutions
of the nonhomogeneous equation
\[
y'+p(x)y=f(x),\tag{A}
\]
then $y_1-y_2$ is a solution of the homogeneous equation
\[
y'+p(x)y=0.\tag{B}
\]

\item %(c)
Use \part{a} to show that if $y_1$ is a solution of
(A) and $y_2$ is a solution of
(B), then $y_1+y_2$ is a solution of
(A).
\end{alist}

\item\label{exer:2.1.47}Some nonlinear equations can be transformed into linear
equations by changing the dependent variable.  Show that if
\[
g'(y)y'+p(x)g(y)=f(x)
\]
where $y$ is a function of $x$ and $g$ is a function of $y$,
then the new dependent variable $z=g(y)$ satisfies the
linear equation
\[
z'+p(x)z=f(x).
\]

\item\label{exer:2.1.48}Solve by the method discussed in \zcref{exer:2.1.47}.

\begin{alist}[columns=2]
% TODO horizontal
\item $(\sec^2y)y'- 3\tan y=-1$
\item $e^{y^2}\left(2yy'+\dfrac{2}{x}\right) =\dfrac{1}{x^2}$
\item $\dfrac{xy'}{y} + 2\ln y=4x^2$
\item $\dfrac{y'}{(1+y)^2} -\dfrac{1}{x(1+y)}=-\dfrac{3}{x^2}$
\end{alist}

\item\label{exer:2.1.49}We've shown that if $p$ and $f$ are continuous on $(a,b)$
then every solution of
\[
y'+p(x)y=f(x)\tag{A}
\]
on $(a,b)$ can be written as $y=uy_1$, where
$y_1$ is a nontrivial solution of the complementary equation for (A)
and $u'=f/y_1$. Now suppose   $f$, $f'$, \dots, $f^{(m)}$ and
$p$, $p'$, \dots, $p^{(m-1)}$ are continuous on $(a,b)$, where $m$
is a positive integer, and define
\begin{eqnarray*}
f_0&=&f,\\ f_j&=&f_{j-1}'+pf_{j-1},\quad 1\le j\le m.
\end{eqnarray*}
Show that
\[
u^{(j+1)}=\frac{f_j}{y_1},\quad 0\le j\le m.
\]

\end{exerciselist}


\section{Separable Equations}\label{sec:separable}

A first order differential equation  is  {\color{blue}\itshape
separable} if it can be written as
\begin{equation} \label{eq:2.2.1}
h(y)y'=g(x),
\end{equation}
where the left side is a product of $y'$ and a function of $y$ and
the right side is a function of $x$. Rewriting a separable
differential equation in this form is called \emph{\color{blue}separation of
variables}. In \zcref{sec:linar_first_order} we used separation of variables
to solve homogeneous linear equations. In this section we'll apply
this method to nonlinear equations.

To see how to solve \eqref{eq:2.2.1}, let's first assume that $y$
is a solution.
Let  $G(x)$ and $H(y)$ be antiderivatives of $g(x)$
and $h(y)$; that is,
\begin{equation}\label{eq:2.2.2}
H'(y)=h(y)\quad\text{and}\quad G'(x)=g(x).
\end{equation}
Then, from the chain rule,
\[
\frac{d}{dx}H(y(x))=H'(y(x))y'(x)=h(y)y'(x).
\]
Therefore \eqref{eq:2.2.1} is equivalent to
\[
\frac{d}{dx}H(y(x))=\frac{d}{dx}G(x).
\]
Integrating both sides of this equation and combining the constants of
integration yields
\begin{equation} \label{eq:2.2.3}
H(y(x))=G(x)+c.
\end{equation}
Although we derived this equation on the assumption that $y$
is a solution of \eqref{eq:2.2.1}, we can now view it differently:
Any differentiable function $y$ that satisfies \eqref{eq:2.2.3}
for some constant $c$ is a solution of \eqref{eq:2.2.1}. To see this, we
differentiate both sides of \eqref{eq:2.2.3}, using the chain rule on the
left, to obtain
\[
H'(y(x))y'(x)=G'(x),
\]
which is equivalent to
\[
h(y(x))y'(x)=g(x)
\]
because of \eqref{eq:2.2.2}.

In conclusion, to solve \eqref{eq:2.2.1} it suffices to find
functions $G=G(x)$ and $H=H(y)$ that satisfy \eqref{eq:2.2.2}.
Then any differentiable function  $y=y(x)$  that satisfies
\eqref{eq:2.2.3} is a solution of \eqref{eq:2.2.1}.

\begin{example}\label{example:2.2.1}
Solve the equation
\[
y'=x(1+y^2).
\]
\end{example}


\solution Separating variables
yields
\[
\frac{y'}{1+y^2}=x.
\]
Integrating yields
\[
\tan^{-1}y=\frac{x^2}{2}+c
\]
Therefore
\[
y=\tan\left(\frac{x^2}{2}+c\right).
\]

\begin{example}\label{example:2.2.2}
\begin{alist}
\item %(a)
Solve the equation
\begin{equation} \label{eq:2.2.4}
y'=-\frac{x}{y}.
\end{equation}

\item %(b)
Solve the initial value problem
\begin{equation} \label{eq:2.2.5}
y'=-\frac{x}{y}, \quad y(1)=1.
\end{equation}

\item %(c)
Solve the initial value problem
\begin{equation} \label{eq:2.2.6}
y'=-\frac{x}{y}, \quad y(1)=-2.
\end{equation}
\end{alist}
\end{example}

\begin{alistSolution}
\item Separating variables
in \eqref{eq:2.2.4}  yields
\[
yy'=-x.
\]
Integrating yields
\[
\frac{y^2}{2}=-\frac{x^2}{2}+c,\mbox{\quad or,\, equivalently, \quad}
x^2+y^2=2c.
\]
The last equation shows that $c$ must be positive if $y$ is to be a
solution of \eqref{eq:2.2.4} on an open interval. Therefore we let
$2c=a^2$ (with $a > 0$) and rewrite the last equation as
\begin{equation} \label{eq:2.2.7}
x^2+y^2=a^2.
\end{equation}
 This equation has two differentiable solutions for $y$ in terms of
$x$:
\begin{equation} \label{eq:2.2.8}
y=\phantom{-} \sqrt{a^2-x^2}, \quad -a < x < a,
\end{equation}
 and
\begin{equation} \label{eq:2.2.9}
y= - \sqrt{a^2-x^2}, \quad -a < x < a.
\end{equation}
The solution curves defined by \eqref{eq:2.2.8}  are semicircles
above the
$x$-axis and  those defined by \eqref{eq:2.2.9}  are semicircles below the
$x$-axis (\zcref{figure:2.2.1}).

\item The solution of  \eqref{eq:2.2.5}
is positive when $x=1$; hence, it  is of the form \eqref{eq:2.2.8}.
 Substituting $x=1$ and $y=1$ into \eqref{eq:2.2.7} to satisfy the
initial condition yields $a^2=2$; hence, the solution of
\eqref{eq:2.2.5}
is
\[
y=\sqrt{2-x^2}, \quad - \sqrt{2}< x < \sqrt{2}.
\]

\item The solution of  \eqref{eq:2.2.6}
is negative when $x=1$ and   is therefore of the form \eqref{eq:2.2.9}.
 Substituting $x=1$ and $y=-2$ into \eqref{eq:2.2.7} to satisfy the
initial condition yields $a^2=5$. Hence, the solution of \eqref{eq:2.2.6}
is
\[
y=- \sqrt{5-x^2}, \quad -\sqrt{5} < x < \sqrt{5}.
\]
\end{alistSolution}

\begin{figure}
  \centering
  \includegraphics[alt={Two semicircles centered at the origin.  One has radius √2 in the upper half; the other has radius √5 in the lower half plane.},
  width=5.67in,height=3.3in,keepaspectratio]{fig020201}
\color{blue}
  \caption{(a) $y=\sqrt{2-x^{2}}$,\, $-\sqrt{2}<x<\sqrt{2}$; \,
(b) $y=-\sqrt{5-x^{2}}$,\, $-\sqrt{5}<x<\sqrt{5}$}
  \label{figure:2.2.1}
\end{figure}


\boxit{Implicit Solutions of Separable Equations}

In \zcref{example:2.2.1,example:2.2.2} we were able to
solve the equation $H(y)=G(x)+c$ to obtain explicit formulas for
solutions of the given separable differential equations. As we'll
see in the next example, this isn't  always possible.
In this situation we must broaden our definition of a solution of a
separable equation. The next theorem provides the basis for this
modification.  We omit the proof, which requires a result from
advanced calculus called as the \emph{\color{blue}implicit function
theorem}.

\begin{theorem}\label{thmtype:2.2.1}
Suppose $g=g(x)$ is continuous on $(a,b)$ and $h=h(y)$
are continuous on $(c,d).$ Let $G$ be an antiderivative of
$g$ on $(a,b)$ and let $H$ be an antiderivative of $h$ on $(c,d).$
Let $x_0$ be an arbitrary point in $(a,b),$ let $y_0$
be a point in $(c,d)$ such that $h(y_0)\ne0,$ and define
\begin{equation} \label{eq:2.2.10}
c=H(y_0)-G(x_0).
\end{equation}
Then there's a function $y=y(x)$ defined on some open interval
$(a_1,b_1),$ where $a\le a_1<x_0<b_1\le b,$ such that  $y(x_0)=y_0$
and
\begin{equation} \label{eq:2.2.11}
H(y)=G(x)+c
\end{equation}
for $a_1<x<b_1$.
Therefore $y$ is a solution of the initial value problem
\begin{equation} \label{eq:2.2.12}
h(y)y'=g(x),\quad y(x_0)=x_0.
\end{equation}
\end{theorem}


It's convenient to say that \eqref{eq:2.2.11} with $c$ arbitrary is an
\emph{\color{blue}implicit solution} of $h(y)y'=g(x)$. Curves defined by
\eqref{eq:2.2.11} are integral curves of $h(y)y'=g(x)$. If $c$ satisfies
\eqref{eq:2.2.10}, we'll say that \eqref{eq:2.2.11} is an {\color{blue}\itshape
implicit
solution of the initial value problem} \eqref{eq:2.2.12}. However, keep
these  points in mind:
\begin{itemize}
\item  For some choices of $c$ there may not be any differentiable
functions $y$ that satisfy \eqref{eq:2.2.11}.
\item  The function $y$ in \eqref{eq:2.2.11} (not \eqref{eq:2.2.11} itself)
is a solution of $h(y)y'=g(x)$.
\end{itemize}

\begin{example}\label{example:2.2.3}
\begin{alist}
\item % (a)
Find implicit solutions of
\begin{equation} \label{eq:2.2.13}
y'=\frac{2x+1}{5y^4+1}.
\end{equation}
\item % (b)
Find an implicit solution of
\begin{equation} \label{eq:2.2.14}
y'=\frac{2x+1}{5y^4+1},\quad y(2)=1.
\end{equation}
\end{alist}
\end{example}

\begin{alistSolution}
\item
 Separating variables yields
\[
(5y^4+1)y'=2x+1.
\]
Integrating yields the implicit solution
\begin{equation} \label{eq:2.2.15}
y^5+y=x^2+x+ c.
\end{equation}
of \eqref{eq:2.2.13}.

\item Imposing the initial condition $y(2)=1$ in
\eqref{eq:2.2.15} yields $1+1=4+2+c$, so $c=-4$. Therefore
\[
y^5+y=x^2+x-4
\]
is an implicit solution of the initial value problem \eqref{eq:2.2.14}.
Although more than one differentiable function $y=y(x)$ satisfies
\eqref{eq:2.2.13}) near $x=2$, it can be shown that there's only
one such function that satisfies the initial condition $y(2)=1$.

\zcref[S]{figure:2.2.2} shows a direction field and some integral curves
for \eqref{eq:2.2.13}.
\end{alistSolution}

\begin{figure}
  \centering
  \includegraphics[alt={A vector field for 1<x<4 and -1<y<2 with arrows pointing up and right.  The arrows are steepest along y=0, and flatten out as they approach y=2.  Several solution curves follow the arrows.},
  width=5.4in,height=3.3in,keepaspectratio]{fig020202}
\color{blue}
  \caption{A direction field and integral curves for
$y'=\dfrac{2x+1}{5y^{4}+1}$}
  \label{figure:2.2.2}
\end{figure}

\boxit{Constant Solutions of Separable Equations}

An equation of the form
\[
y'=g(x)p(y)
\]
 is separable, since it can be rewritten as
\[
\frac{1}{p(y)}y'=g(x).
\]
 However, the division by  $p(y)$
is not legitimate  if $p(y)=0$ for some values of $y$.  The
next two examples show how to deal with this problem.

\begin{example}\label{example:2.2.4}
Find all solutions of
\begin{equation} \label{eq:2.2.16}
y'=2xy^2.
\end{equation}
\end{example}
\solution Here we must divide by $p(y)=y^2$ to separate variables.
This isn't  legitimate if $y$ is a solution of \eqref{eq:2.2.16} that
equals zero for some value of $x$. One such solution can be found by
inspection:  $y \equiv 0$. Now suppose   $y$ is a solution
of \eqref{eq:2.2.16} that isn't  identically zero. Since $y$ is continuous
there must be an interval on which $y$ is never zero. Since division
by $y^2$ is legitimate for $x$ in this interval, we can separate
variables in \eqref{eq:2.2.16} to obtain
\[
\frac{y'}{y^2}=2x.
\]
 Integrating this yields
\[
-\frac{1}{y}=x^2+c,
\]
which is equivalent to
\begin{equation} \label{eq:2.2.17}
y=-\frac{1}{x^2+c}.
\end{equation}

We've now shown that if $y$ is a solution of \eqref{eq:2.2.16} that is
not identically zero, then $y$ must be of the form \eqref{eq:2.2.17}. By
substituting \eqref{eq:2.2.17} into \eqref{eq:2.2.16}, you can verify that
\eqref{eq:2.2.17} is  a solution of \eqref{eq:2.2.16}. Thus,
 solutions of \eqref{eq:2.2.16} are $y\equiv0$ and the functions
of the form \eqref{eq:2.2.17}. Note that the solution $y\equiv0$ isn't  of
the form \eqref{eq:2.2.17} for any value of $c$.

\zcref[S]{figure:2.2.3} shows a direction field and some integral
curves for \eqref{eq:2.2.16}

\begin{figure}
  \centering
  \includegraphics[alt={A vector field for -2<x,y<2, with arrows pointing to the right and downward when x<0 and right and upward when x>0.  For y<0, solution curves start along y=0, curve downward to bottom out when x=0, and then increase back to finish along y=0.  For y>0, solution curves start at vertical asymptotes, bottom out when x=0, and then increase to finish at other vertical asymptotes.},
  width=5.4in,height=3.3in,keepaspectratio]{fig020203}
\color{blue}
  \caption{A direction field and  integral curves for $y'=2xy^{2}$}
  \label{figure:2.2.3}
\end{figure}

\begin{example}\label{example:2.2.5}
Find all solutions of
\begin{equation} \label{eq:2.2.18}
y'=\frac{1}{2}x(1-y^2).
\end{equation}
\end{example}

\solution Here we must divide by $p(y)=1-y^2$ to separate variables.
This isn't  legitimate if $y$ is a solution of \eqref{eq:2.2.18} that
equals $\pm1$ for some value of $x$. Two such solutions can be found
by inspection: $y \equiv 1$ and $y\equiv-1$. Now suppose
$y$ is a solution of \eqref{eq:2.2.18} such that $1-y^2$ isn't
identically zero. Since $1-y^2$ is continuous there must be an
interval on which $1-y^2$ is never zero. Since division by $1-y^2$ is
legitimate for $x$ in this interval, we can separate variables in
\eqref{eq:2.2.18} to obtain
\[
\frac{2y'}{y^2-1}=-x.
\]
 A partial fraction expansion on the left yields
\[
\left[\frac{1}{y-1}-\frac{1}{y+1}\right]y'=-x,
\]
and integrating  yields
\[
\ln\left|\frac{y-1}{y+1}\right|=-\frac{x^2}{2}+k;
\]
 hence,
\[
\left|\frac{y-1}{y+1}\right|=e^ke^{-x^2/2}.
\]
 Since $y(x)\ne\pm1$ for $x$ on the interval under discussion, the
quantity
$(y-1)/(y+1)$  can't  change sign in this interval. Therefore
we can rewrite the last equation as
\[
\frac{y-1}{y+1}=ce^{-x^2/2},
\]
  where $c=\pm e^k$, depending upon the sign of $(y-1)/(y+1)$ on the
interval.
 Solving for  $y$ yields
\begin{equation} \label{eq:2.2.19}
y=\frac{1+ce^{-x^2/2}}{1-ce^{-x^2/2}}.
\end{equation}

We've now shown that if $y$ is a solution of \eqref{eq:2.2.18} that  is
not identically equal to $\pm1$, then $y$ must be as in
\eqref{eq:2.2.19}. By substituting \eqref{eq:2.2.19} into \eqref{eq:2.2.18} you
 can verify that \eqref{eq:2.2.19} is  a solution of
\eqref{eq:2.2.18}. Thus,  the solutions of \eqref{eq:2.2.18}
are $y\equiv1$, $y\equiv-1$ and the functions of the form
\eqref{eq:2.2.19}. Note that the constant solution $y \equiv 1$ can be
obtained from this formula by taking $c=0$;     however, the other
constant solution, $y \equiv -1$, can't be obtained in this way.

\zcref[S]{figure:2.2.4} shows a direction field and some integral curves for
\eqref{eq:2.2.18}.

\begin{figure}
  \centering
  \includegraphics[alt={A vector field for -2<x<2 and -3<y<3 with trivial solutions y=±1.  When y>1, the curves increase to a peak at x=0 before decreasing to y=1.  When -1<y<1, the curves decrease and are flat near x=0 before increasing to y=1.  When y<-1, the curves begin at a vertical asymptote, increase to a maximum at x=0, and then go to another vertical asymptote.},
  width=5.67in,height=3.3in,keepaspectratio]{fig020204}
\color{blue}
\caption{
A direction field and integral curves for
$y'=\dfrac{x(1-y^2)}{2}$}
  \label{figure:2.2.4}
\end{figure}

\boxit{Differences Between Linear and Nonlinear Equations}

\zcref[S]{thmtype:2.1.2} states that if $p$ and $f$ are continuous on
$(a,b)$ then every solution of
\[
y'+p(x)y=f(x)
\]
 on  $(a,b)$  can be obtained by choosing a value for the constant
$c$  in the general solution, and   if
$x_0$ is any point in  $(a,b)$ and $y_0$ is arbitrary, then the
initial value problem
\[
y'+p(x)y=f(x),\quad y(x_0)=y_0
\]
 has a solution on  $(a,b)$.

The not true for nonlinear equations. First, we
saw in \zcref{example:2.2.4,example:2.2.5} that a
nonlinear
equation may have solutions that can't be obtained by choosing a
specific value of a constant appearing in a one-parameter family of
solutions. Second, it is in general impossible to determine the
interval of validity of a solution to an initial value problem for a
nonlinear equation by simply examining the equation, since the
interval of validity may depend on the initial condition. For
instance, in \zcref{example:2.2.2} we saw that the solution of
\[
\frac{dy}{dx}=-\frac{x}{y},\quad y(x_0)=y_0
\]
is valid on $(-a,a)$, where $a=\sqrt{x_0^2+y_0^2}$.

\begin{example}\label{example:2.2.6}
Solve the initial value problem
\[
y'=2xy^2, \quad y(0)=y_0
\]
 and determine the  interval of validity of  the solution.
\end{example}


\solution  First suppose   $y_0\ne0$. From
\zcref{example:2.2.4}, we know that $y$ must be of the form
\begin{equation} \label{eq:2.2.20}
y=-\frac{1}{x^2+c}.
\end{equation}
 Imposing the initial condition shows  that $c=-1/y_0$.
Substituting this into \eqref{eq:2.2.20} and rearranging terms yields
the solution
\[
y= \frac{y_0}{1-y_0x^2}.
\]
This is also the solution if $y_0=0$. If $y_0<0$,
the denominator isn't  zero for any value of $x$, so the
the solution is valid on $(-\infty,\infty)$. If $y_0>0$, the solution
 is  valid only on $(-1/\sqrt{y_0},1/\sqrt{y_0})$.

\exercises

\emph{In \zcref[range]{exer:2.2.1,exer:2.2.6} find all solutions.}

\begin{exerciselist}
\item(1)\label{exer:2.2.1}$y'=\dfrac{3x^2+2x+1}{y-2}$
\item(1)\label{exer:2.2.2}$(\sin x)(\sin y)+(\cos y)y'=0$
\item(1)\label{exer:2.2.3}$xy'+y^2+y=0$
\item(1)\label{exer:2.2.4}$y' \ln |y|+x^2y= 0$

\item\label{exer:2.2.5}$(3y^3+3y \cos y+1)y'+\dfrac{(2x+1)y}{1+x^2}=0$


\item\label{exer:2.2.6}$x^2yy'=(y^2-1)^{3/2}$

\exercisetext{In \zcref[range]{exer:2.2.7,exer:2.2.10} find all
solutions. Also, plot a direction field and some integral curves on
the indicated rectangular region.}


\item\label{exer:2.2.7}\CGex $y'=x^2(1+y^2);   \; \{-1\le x\le1,\
-1\le y\le1\}$


 \item\label{exer:2.2.8}\CGex$ y'(1+x^2)+xy=0 ;  \; \{-2\le x\le2,\ -1\le
y\le1\}$

\item\label{exer:2.2.9}\CGex $y'=(x-1)(y-1)(y-2);   \; \{-2\le x\le2,\
-3\le y\le3\}$


 \item\label{exer:2.2.10}\CGex $(y-1)^2y'=2x+3;   \; \{-2\le x\le2,\ -2\le
y\le5\}$

\exercisetext{In \zcref{exer:2.2.11,exer:2.2.12} solve the initial value problem.}

\item\label{exer:2.2.11}$y'=\dfrac{x^2+3x+2}{y-2}, \quad y(1)=4$

\item\label{exer:2.2.12}$y'+x(y^2+y)=0, \quad y(2)=1$

\exercisetext{In \zcref[range]{exer:2.2.13,exer:2.2.16} solve
the initial value problem and graph the solution.}

\item\label{exer:2.2.13}\CGex $(3y^2+4y)y'+2x+\cos x=0, \quad y(0)=1$



\item\label{exer:2.2.14}\CGex
$y'+\dfrac{(y+1)(y-1)(y-2)}{x+1}=0, \quad y(1)=0$

\item\label{exer:2.2.15}\CGex $y'+2x(y+1)=0, \quad y(0)=2$


\item\label{exer:2.2.16}\CGex $y'=2xy(1+y^2),\quad y(0)=1$

\exercisetext{In \zcref[range]{exer:2.2.17,exer:2.2.23} solve the initial
value problem and find the interval of validity of the solution.}

\item\label{exer:2.2.17}$y'(x^2+2)+ 4x(y^2+2y+1)=0, \quad y(1)=-1$

\item\label{exer:2.2.18}$y'=-2x(y^2-3y+2), \quad y(0)=3$

\item(1)\label{exer:2.2.19}$y'=\dfrac{2x}{1+2y},\quad y(2)=0$
\item(1)\label{exer:2.2.20}$y'=2y-y^2,\quad y(0)=1$

\item\label{exer:2.2.21}$x+yy'=0, \quad y(3) =-4$

\item\label{exer:2.2.22}$y'+x^2(y+1)(y-2)^2=0, \quad y(4)=2$

\item\label{exer:2.2.23}$(x+1)(x-2)y'+y=0, \quad y(1)=-3$

\item\label{exer:2.2.24}Solve $y'=\dfrac{(1+y^2)}{(1+x^2)}$ explicitly.
\hint{Use the identity
$\tan(A+B)=\dfrac{\tan A+\tan B}{1-\tan A\tan B}$.}

\item\label{exer:2.2.25}Solve $y'\sqrt{1-x^2}+\sqrt{1-y^2}=0$ explicitly.
\hint{Use the identity
$\sin(A-B)=\sin A\cos B-\cos A\sin B$.}

\item\label{exer:2.2.26}Solve $y'=\dfrac{\cos x}{\sin y},\quad y (\pi)=\dfrac{\pi}{2}$
explicitly. \hint{Use the identity $\cos(x+\pi/2)=-\sin x$ and the
periodicity of the cosine.}

\item\label{exer:2.2.27}Solve the initial value problem
\[
y'=ay-by^2,\quad y(0)=y_0.
\]
Discuss the behavior  of the solution if
\part{a} $y_0\ge0$; \part{b} $y_0<0$.

\item\label{exer:2.2.28}The population $P=P(t)$ of a  species satisfies
the logistic equation
\[
P'=aP(1-\alpha P)
\]
 and  $P(0)=P_0>0$.  Find $P$  for $t>0$, and  find
$\lim_{t\to\infty}P(t)$.



\item\label{exer:2.2.29}An epidemic spreads through a population at a rate
proportional to the product of the number of people already infected
and the number of people susceptible, but not yet infected. Therefore,
 if $S$ denotes the total population of susceptible people and
$I=I(t)$ denotes the number of infected people at time $t$, then
\[
I'=rI(S-I),
\]
where $r$ is a positive constant. Assuming that $I(0)=I_0$,
find $I(t)$ for $t>0$, and show that
$\lim_{t\to\infty}I(t)=S$.

\item\label{exer:2.2.30}\Lex
The result of \zcref{exer:2.2.29} is discouraging: if any
susceptible member of the group is
initially infected, then in the long run all susceptible members are
infected!
 On a more hopeful note, suppose   the disease spreads
according to the model of \zcref{exer:2.2.29}, but there's a
medication that cures the infected population at a rate proportional
to the number of infected individuals. Now the equation for the number
of infected individuals becomes
\[
I'=rI(S-I)-qI\tag{A}
\]
where $q$ is a positive constant.
\begin{alist}
\item % (a)
Choose  $r$ and $S$ positive.
By plotting  direction fields and solutions of (A) on
suitable rectangular grids
\[
R=\{0\le t \le T,\ 0\le I \le d\}
\]
in the $(t,I)$-plane, verify that if $I$ is any solution of
(A) such that $I(0)>0$, then $\lim_{t\to\infty}I(t)=S-q/r$
if $q<rS$ and $\lim_{t\to\infty}I(t)=0$ if $q\ge rS$.
\item % (b)
To verify the experimental results of \part{a}, use separation of
variables to solve (A) with initial condition $I(0)=I_0>0$, and find
$\lim_{t\to\infty}I(t)$. \hint{There are three cases to consider:
\part{i} $q<rS$;     \part{ii}   $q>rS$;   \part{iii} $q=rS$.}
\end{alist}

\item\label{exer:2.2.31}\Lex
Consider the differential equation
\[
y'=ay-by^2-q,\tag{A}
\]
where $a$, $b$ are positive constants, and $q$ is an arbitrary
constant. Suppose  $y$ denotes a solution of this equation
that satisfies the initial condition $y(0)=y_0$.

\begin{alist}
\item % (a)
Choose  $a$ and $b$ positive and  $q<a^2/4b$.
By plotting  direction fields and solutions of (A) on
suitable rectangular grids
\[
R=\{0\le t \le T,\ c\le y \le d\}\tag{B}
\]
in the $(t,y)$-plane,  discover that there are numbers $y_1$ and
$y_2$ with $y_1<y_2$ such that if $y_0>y_1$ then
$\lim_{t\to\infty}y(t)=y_2$, and if $y_0<y_1$ then $y(t)=-\infty$
for some finite value of~$t$. (What happens if $y_0=y_1$?)



\item % (b)
Choose $a$ and $b$ positive and $q=a^2/4b$. By plotting direction
fields
and solutions of (A) on suitable rectangular grids of the form (B),
discover that there's a number $y_1$ such that if $y_0\ge y_1$ then
$\lim_{t\to\infty}y(t)=y_1$, while if $y_0<y_1$ then  $y(t)=-\infty$
for some finite value of $t$.

\item % (c)
Choose positive $a$, $b$ and $q>a^2/4b$. By plotting direction fields
and solutions of (A) on suitable rectangular grids of the form (B),
discover that
no matter what $y_0$ is, $y(t)=-\infty$ for some finite value of $t$.



\item % (d)
Verify your results experiments analytically.
Start by separating variables in (A)  to obtain
\[
\frac{y'}{ay-by^2-q}=1.
\]
To decide what to do next you'll have to use the
quadratic formula. This should lead you to see why there are
three cases.  Take it from there!

Because of its role in the transition between these
three cases,  $q_0=a^2/4b$ is called a \emph{\color{blue}bifurcation value}
of $q$. In general, if $q$ is a parameter in any differential
equation,  $q_0$ is said to be a bifurcation value of
$q$ if the nature of the solutions of the equation with $q<q_0$
is qualitatively different from the nature of the solutions
with $q>q_0$.
\end{alist}

\item\label{exer:2.2.32}\Lex
By plotting direction fields and solutions of
\[
y'=qy-y^3,
\]
convince yourself that $q_0=0$ is a bifurcation value of $q$
for this equation. Explain what makes you draw this conclusion.

\item\label{exer:2.2.33}Suppose a disease spreads according to the model of
\zcref[S]{exer:2.2.29}, but
there's a medication that cures the infected population at a constant
rate of $q$ individuals per unit  time, where $q>0$.
Then
the equation for the number of infected individuals becomes
\[
I'=rI(S-I)-q.
\]
Assuming that $I(0)=I_0>0$, use the results of
\zcref{exer:2.2.31} to describe what happens as $t\to\infty$.

\item\label{exer:2.2.34}Assuming that $p \not\equiv 0$, state conditions under which the linear
equation
\[
y'+p(x)y=f(x)
\]
is separable.  If the equation satisfies these conditions, solve it by
separation of variables and by the method developed in
\zcref{sec:linar_first_order}.

\exercisetext{Solve the equations in
\zcref[range]{exer:2.2.35,exer:2.2.38} using variation of
parameters followed by separation of variables.}

\item(1)\label{exer:2.2.35}$y'+y=\dfrac{2xe^{-x}}{1+ye^x}$
\item(1)\label{exer:2.2.36}$xy'-2y=\dfrac{x^6}{y+x^2}$
\item(1)\label{exer:2.2.37}$y'-y=\dfrac{(x+1)e^{4x}}{(y+e^x)^2}$
\item(1)\label{exer:2.2.38}$y'-2y=\dfrac{xe^{2x}}{1-ye^{-2x}}$


\item\label{exer:2.2.39}Use variation of parameters to show that the solutions of the
following equations are of the form $y=uy_1$, where $u$ satisfies a
separable equation $u'=g(x)p(u)$. Find $y_1$ and $g$ for each
equation.

\begin{alist}[columns=2]
% TODO horizontal
\item $xy'+y=h(x)p(xy)$
\item $xy'-y=h(x) p\left(\dfrac{y}{x}\right)$
\item $y'+y=h(x) p(e^xy)$
\item $xy'+ry=h(x) p(x^ry)$
\item $y'+\dfrac{v'(x)}{v(x)}y= h(x)p\left(v(x)y\right)$
\end{alist}
\end{exerciselist}


\section{Existence and Uniqueness of Solutions of Nonlinear Equations}
\label{sec:exist_unique}

Although  there are  methods for
 solving some nonlinear equations, it's
impossible to find  useful formulas for the solutions of most.
Whether we're looking for  exact solutions or numerical
approximations, it's useful to know  conditions that imply the
existence and uniqueness of solutions of initial value problems for
nonlinear equations. In this section we state  such a condition and
illustrate it with examples.

\begin{figure}
  \centering
  \includegraphics[alt={The rectangle a<x<b and c<y<d with dots insides and a dashed boundary.},
  trim=0 50 0 60,clip,width=5.67in,height=3.3in,keepaspectratio]{fig020301}
\color{blue}
  \caption{An open rectangle}
  \label{figure:2.3.1}
\end{figure}


Some terminology:
 an \emph{\color{blue}open rectangle $R$}
 is a set of points $(x,y)$ such that
\[
a<x<b\mbox{\quad  and \quad}c<y<d
\]
 (\zcref{figure:2.3.1}).  We'll denote this set by
$R:  \{ a < x < b, c < y < d \}$.
 ``Open'' means that the
boundary rectangle (indicated by the dashed lines in
\zcref{figure:2.3.1}) isn't  included in  $R$ .

The next theorem gives sufficient conditions for existence and
uniqueness of solutions of initial value problems for first order
nonlinear differential equations. We omit the proof, which is beyond
the scope of this book.

\begin{theorem}\label{thmtype:2.3.1}
\begin{alist}
\item %(a)
 If  $f$ is continuous
on an open rectangle
\[
R:  \{ a < x < b, c < y < d \}
\]
 that contains $(x_0,y_0)$
then  the initial value problem
\begin{equation} \label{eq:2.3.1}
y'=f(x,y), \quad y(x_0)=y_0
\end{equation}
 has at least one solution  on some open subinterval
of  $(a,b)$ that contains $x_0.$

\item % (b)
 If  both $f$ and  $f_y$ are
 continuous on $R$ then \eqref{eq:2.3.1} has a unique
solution on some open subinterval  of $(a,b)$ that contains $x_0$.
\end{alist}
\end{theorem}


It's important to understand exactly what \zcref{thmtype:2.3.1}
says.

\begin{itemize}
\item (a) is an \emph{\color{blue}existence theorem}. It guarantees that a
solution exists on some open interval that contains $x_0$, but provides
no information on how to find the solution, or to determine the open
interval on which it exists. Moreover, \part{a} provides no
information on the number of solutions that \eqref{eq:2.3.1} may have. It
leaves open the possibility that \eqref{eq:2.3.1} may have two or more
solutions that differ for values of $x$ arbitrarily close to $x_0$. We
will see in \zcref{example:2.3.6} that this can happen.

\item (b) is a \emph{\color{blue}uniqueness theorem}. It guarantees that
\eqref{eq:2.3.1} has a unique solution on some open interval (a,b) that
contains $x_0$. However, if $(a,b)\ne(-\infty,\infty)$,
\eqref{eq:2.3.1} may have more than one solution on a larger interval
that contains $(a,b)$. For example, it may happen that $b<\infty$ and all
solutions have the same values on $(a,b)$, but two solutions $y_1$ and
$y_2$ are defined on some interval $(a,b_1)$ with $b_1>b$, and have
different values for $b<x<b_1$;   thus, the graphs of the $y_1$ and
$y_2$ ``branch off'' in different directions at $x=b$. (See
\zcref{example:2.3.7} and \zcref{figure:2.3.3}). In this case,
continuity implies that $y_1(b)=y_2(b)$ (call their common value
$\overline y$), and  $y_1$ and $y_2$ are both solutions of the initial
value problem
\begin{equation} \label{eq:2.3.2}
y'=f(x,y),\quad y(b)=\overline y
\end{equation}
that differ on every open interval that contains $b$. Therefore
$f$ or $f_y$ must have a discontinuity at some point in each open
rectangle that contains $(b,\overline y)$, since if this were not so,
\eqref{eq:2.3.2} would have a unique solution on some open interval
that contains $b$. We leave it to you to give a similar analysis of the
case where $a>-\infty$.
 \end{itemize}

\begin{example}\label{example:2.3.1}
Consider  the initial value problem
\begin{equation} \label{eq:2.3.3}
y'=\frac{x^2-y^2 }{1+x^2+y^2}, \quad y(x_0)=y_0.
\end{equation}
Since
\[
f(x,y)  = \frac{x^2-y^2 }{1+x^2+y^2}\quad\text{and}\quad 
f_y(x,y)  =  -\frac{2y(1+2x^2)}{(1+x^2+y^2)^2}
\]
are continuous for all $(x,y)$, \zcref{thmtype:2.3.1} implies that
if $(x_0,y_0)$ is arbitrary, then
\eqref{eq:2.3.3} has a unique solution on some open interval that contains
$x_0$.
\end{example}


\begin{example}\label{example:2.3.2}
Consider  the initial value problem
\begin{equation} \label{eq:2.3.4}
y'=\frac{x^2-y^2 }{x^2+y^2}, \quad y(x_0)=y_0.
\end{equation}
Here
\[
f(x,y)  =  \frac{x^2-y^2 }{x^2+y^2}\quad\text{and}\quad 
f_y(x,y)  =   -\frac{4x^2y }{(x^2+y^2)^2}
\]
are continuous everywhere except at $(0,0)$. If $(x_0,y_0)
\ne(0,0)$, there's an open rectangle $R$ that contains
$(x_0,y_0)$ that does not contain $(0,0)$. Since $f$ and $f_y$ are
continuous on $R$, \zcref{thmtype:2.3.1} implies that if
$(x_0,y_0)\ne(0,0)$ then
\eqref{eq:2.3.4}
has a unique solution on some open interval that contains $x_0$.
\end{example}

\begin{example}\label{example:2.3.3}
Consider the initial value problem
\begin{equation} \label{eq:2.3.5}
y'=\frac{x+y}{x-y},\quad y(x_0)=y_0.
\end{equation}
Here
\[
f(x,y)  =  \frac{x+y}{x-y}\quad\text{and}\quad
f_y(x,y)  =  \frac{2x}{(x-y)^2}
\]
are continuous everywhere except on the line $y=x$. If $y_0\ne x_0$,
there's an open rectangle $R$ that contains $(x_0,y_0)$ that
does not intersect the line $y=x$. Since $f$ and $f_y$ are continuous
on $R$, \zcref{thmtype:2.3.1} implies that if $y_0\ne x_0$,
\eqref{eq:2.3.5} has a unique solution on some open interval that contains
$x_0$.
\end{example}

\begin{example}\label{example:2.3.4}
In \zcref{example:2.2.4} we saw that the
solutions of
\begin{equation} \label{eq:2.3.6}
y'=2xy^2
\end{equation}
are
 \[
y\equiv0\quad\text{and}\quad  y=-\frac{1 }{x^2+c},
\]
where $c$ is an arbitrary constant. In particular, this implies that
no solution of \eqref{eq:2.3.6} other than $y\equiv0$ can equal zero for
any value of $x$. Show that \zcref{thmtype:2.3.1}\part{b} implies
this.
\end{example}

\solution
We'll obtain a contradiction
by assuming that \eqref{eq:2.3.6} has a solution $y_1$ that equals
zero for some value of $x$, but isn't  identically zero. If $y_1$
has this property,  there's a point $x_0$ such that $y_1(x_0)=0$,
but $y_1(x)\ne0$ for some value of $x$ in every open interval that contains
$x_0$. This means that the initial value problem
\begin{equation} \label{eq:2.3.7}
y'=2xy^2,\quad y(x_0)=0
\end{equation}
has two solutions $y\equiv0$ and $y=y_1$ that differ for some value of
$x$ on every open interval that contains $x_0$. This contradicts
\zcref{thmtype:2.3.1}(b), since in \eqref{eq:2.3.6} the functions
\[
f(x,y)=2xy^2 \quad\text{and}\quad  f_y(x,y)= 4xy.
\]
are both continuous
for all $(x,y)$, which implies that \eqref{eq:2.3.7} has a unique
solution on some open interval that contains $x_0$.

\begin{example}\label{example:2.3.5}
Consider the initial value problem
\begin{equation} \label{eq:2.3.8}
y' = \frac{10}{3}xy^{2/5}, \quad y(x_0) = y_0.
\end{equation}
\begin{alist}
\item %(a)
For what points $(x_0,y_0)$ does \zcref{thmtype:2.3.1}\part{a}
imply that
\eqref{eq:2.3.8} has a solution?

\item %(b)
For what points $(x_0,y_0)$ does \zcref{thmtype:2.3.1}\part{b}
imply that
\eqref{eq:2.3.8} has a  unique solution on some open interval that contains
$x_0$?
\end{alist}
\end{example}

\begin{alistSolution}
\item Since
\[
f(x,y) = \frac{10}{3}xy^{2/5}
\]
is continuous for all $(x,y)$,  \zcref{thmtype:2.3.1}
implies that \eqref{eq:2.3.8} has a solution for every $(x_0,y_0)$.

\item Here
\[
f_y(x,y) = \frac{4 }{3}xy^{-3/5}
\]
is continuous for all $(x,y)$ with $y\ne 0$. Therefore, if $y_0\ne0$
there's an open rectangle on which both $f$ and $f_y$ are
continuous, and \zcref{thmtype:2.3.1} implies that \eqref{eq:2.3.8} has
a unique solution on some open interval that contains $x_0$.

If $y=0$ then $f_y(x,y)$ is undefined, and therefore discontinuous;
hence, \zcref{thmtype:2.3.1} does not apply to \eqref{eq:2.3.8} if
$y_0=0$.
\end{alistSolution}

\begin{example}\label{example:2.3.6}
\zcref[S]{example:2.3.5} leaves open the possibility  that the
initial value problem
\begin{equation} \label{eq:2.3.9}
y'=\frac{10 }{3}xy^{2/5}, \quad y(0)=0
\end{equation}
has more than one solution on every open interval that contains $x_0=0$.  Show
that this is true.
\end{example}

\solution
By inspection, $y\equiv0$ is a solution of the
differential equation
\begin{equation} \label{eq:2.3.10}
y'=\frac{10 }{3} xy ^{2/5}.
\end{equation}
 Since $y\equiv0$ satisfies the initial condition $y(0)=0$,
it's a solution of \eqref{eq:2.3.9}.

Now suppose  $y$ is a solution of \eqref{eq:2.3.10} that isn't
identically zero.
 Separating variables in \eqref{eq:2.3.10} yields
\[
 y^{-2/5}y'=\frac{10 }{3}x
\]
on any open interval where $y$ has no zeros.
 Integrating this and rewriting the arbitrary constant as $5c/3$ yields
\[
\frac{5}{3}y^{3/5} = \frac{5}{3}(x^2+c).
\]
Therefore
\begin{equation} \label{eq:2.3.11}
y = (x^2+c)^{5/3}.
\end{equation}

Since we divided by $y$ to separate variables in \eqref{eq:2.3.10}, our
derivation of \eqref{eq:2.3.11} is legitimate only on open intervals where $y$
has no zeros. However, \eqref{eq:2.3.11} actually defines $y$ for all $x$,
and differentiating \eqref{eq:2.3.11} shows that
\[
y'=\frac{10 }{3}x(x^2+c)^{2/3}=\frac{10 }{3}xy^{2/5},\quad-\infty<x<\infty.
\]
Therefore \eqref{eq:2.3.11} satisfies \eqref{eq:2.3.10} on
$(-\infty,\infty)$
even if $c\le 0$, so that $y(\sqrt{|c|})=y(-\sqrt{|c|})=0$. In
particular, taking $c=0$ in \eqref{eq:2.3.11} yields
\[
y=x^{10/3}
\]
as a second solution of \eqref{eq:2.3.9}.
 Both solutions are defined on
$(-\infty,\infty)$, and they differ on every
open interval that contains $x_0=0$ (see \zcref{figure:2.3.2}.)
In fact, there are \emph{\color{blue}four} distinct solutions of
\eqref{eq:2.3.9}  defined on $(-\infty,\infty)$ that differ
from each other on every open interval that contains $x_0=0$.
  Can you identify the other two?

\begin{figure}
  \centering
%  \includegraphics[%bb=-78 148 689 643,
%  width=5.67in,height=3.3in,keepaspectratio]{fig020302}
\begin{tikzpicture}[alt={A curve that begins near (-2,-10), crosses horizontally through the origin, and finishes in a mirror image near (2,10).  A second line follows the x axis.}]
\begin{axis}[width=\textwidth, height=3.3in,
axis y line=middle,axis x line=middle, tick label style={font=\footnotesize},
ymin=-10,ymax=10,xmin=-2,xmax=2,
%xtick=\empty,ytick={5},yticklabels={$T_0$},
name=myplot] %
\addplot[domain=-2:0,draw=blue]{-(-x)^(10/3)};
\addplot[domain=0:2,draw=blue]{x^(10/3)};
\addplot[domain=-2:2,draw=blue]{0};
%\pgfplotsinvokeforeach{1,2,...,5}{
%\addplot[draw=blue, domain=0:2, smooth] {5+#1*exp(-x)};
%\addplot[draw=blue, domain=0:2, smooth] {5-#1*exp(-x)};
%}
\end{axis}
\node [right] at (myplot.right of origin) {\scriptsize $x$};
\node [above] at (myplot.above origin) {\scriptsize $y$};
\end{tikzpicture}
\color{blue}
\caption{Two solutions ($y=0$ and $y=x^{10/3}$) of \eqref{eq:2.3.9}  that
differ on every interval  containing $x_{0}=0$}
  \label{figure:2.3.2}
\end{figure}

\begin{example} \label{example:2.3.7}
From \zcref{example:2.3.5}, the initial value
problem
\begin{equation} \label{eq:2.3.12}
y'=\frac{10 }{3}xy^{2/5}, \quad y(0)=-1
\end{equation}
 has a  unique solution on some open interval that contains $x_0=0$.
Find a solution and determine the largest open interval $(a,b)$ on
which it's unique.
\end{example}

\solution
Let $y$ be any solution of \eqref{eq:2.3.12}. Because of the initial
condition $y(0)=-1$ and the continuity of $y$, there's an open interval
$I$ that contains $x_0=0$ on which $y$ has no zeros, and is consequently
of the form \eqref{eq:2.3.11}. Setting $x=0$ and $y=-1$ in \eqref{eq:2.3.11}
yields $c=-1$, so
\begin{equation} \label{eq:2.3.13}
y=(x^2-1)^{5/3}
\end{equation}
for $x$  in $I$.
Therefore every solution of \eqref{eq:2.3.12} differs from zero
and is given by \eqref{eq:2.3.13} on $(-1,1)$; that is,
\eqref{eq:2.3.13} is the unique solution of \eqref{eq:2.3.12} on
$(-1,1)$.
 This is the largest open interval on which
\eqref{eq:2.3.12} has a unique solution. To see this, note that
\eqref{eq:2.3.13} is a solution of \eqref{eq:2.3.12} on $(-\infty,\infty)$.
From {Exercise~2.2.\ref{exer:2.2.15}}, there are infinitely many
other solutions
of \eqref{eq:2.3.12} that differ from \eqref{eq:2.3.13} on every open interval
larger than
$(-1,1)$. One such solution is
\[
y =
\begin{cases}
(x^2-1)^{5/3}, & -1 \le x \le 1, \\
0,             & |x|>1.
\end{cases}
\]
(see \zcref{figure:2.3.3}).

\noindent
\begin{minipage}[t]{0.48\linewidth}
\tagstructbegin{tag=Sect}
\includegraphics[alt={A blue curve beginning near (-2,6), decreasing to horizontally cross the x axis at x=1, bottoming out at (0,-1), and finishing in a mirror image near (2,6).  A second blue curve goes from (-2,0) to (-1,0) and from (1,0) to (2,0).},
width=\linewidth,height=2.44in,keepaspectratio]{fig020303}
\captionof{figure}{Two solutions of \eqref{eq:2.3.12} on $(-\infty,\infty)$
that coincide on $(-1,1)$, but on no larger open interval}
\label{figure:2.3.3}
\tagstructend
\end{minipage}%
\hfill
\begin{minipage}[t]{0.48\linewidth}
\tagstructbegin{tag=Sect}\includegraphics[alt={A blue curve beginning near (-2,6), decreasing to bottom out at (0,1), and finishing in a mirror image near (2,6).},
width=\linewidth,height=2.44in,keepaspectratio]{fig020304}
\captionof{figure}{The unique solution of \eqref{eq:2.3.14}}
\label{figure:2.3.4}
\tagstructend
\end{minipage}


\begin{example}\label{example:2.3.8}
From \zcref{example:2.3.5}, the initial value
problem
\begin{equation} \label{eq:2.3.14}
y'=\frac{10 }{3}xy^{2/5}, \quad y(0)=1
\end{equation}
 has a unique solution on some open interval that contains $x_0=0$.
Find the solution and determine the largest open interval on which it's
unique.
\end{example}

\solution
Let $y$ be any solution of \eqref{eq:2.3.14}. Because of the initial
condition $y(0)=1$ and the continuity of $y$, there's an open interval
$I$ that contains $x_0=0$ on which $y$ has no zeros, and is consequently
of the form \eqref{eq:2.3.11}. Setting $x=0$ and $y=1$ in \eqref{eq:2.3.11}
yields $c=1$, so
\begin{equation} \label{eq:2.3.15}
y=(x^2+1)^{5/3}
\end{equation}
for $x$ in $I$. Therefore every solution of \eqref{eq:2.3.14}
differs from zero and is given by \eqref{eq:2.3.15} on $(-\infty,\infty)$;
that is, \eqref{eq:2.3.15} is the unique solution of \eqref{eq:2.3.14} on
$(-\infty,\infty)$.
\zcref[S]{figure:2.3.4} shows the graph of this solution.

\exercises

\noindent
\emph{In \zcref[range]{exer:2.3.1,exer:2.3.13} find all $(x_0,y_0)$ for
which \zcref{thmtype:2.3.1}
implies that the initial value problem $y'=f(x,y),\  y(x_0)=y_0$ has
(a) a solution and (b) a  unique solution on some open interval that
contains $x_0$.}

\begin{exerciselist}

\item(1)\label{exer:2.3.1}$y'=\dfrac{x^2+y^2}{\sin x}$
\item(1)\label{exer:2.3.2}$y'=\dfrac{e^x+y }{x^2+y^2}$
\item(1)\label{exer:2.3.3}$y'= \tan xy$
\item(1)\label{exer:2.3.4}$y'=\dfrac{x^2+y^2}{\ln xy}$
\item(1)\label{exer:2.3.5}$y'=(x^2+y^2)y^{1/3}$
\item(1)\label{exer:2.3.6}\ $y'=2xy$
\item(1)\label{exer:2.3.7}$y'=\ln(1+x^2+y^2)$
\item(1)\label{exer:2.3.8}$y'=\dfrac{2x+3y}{x-4y}$
\item(1)\label{exer:2.3.9}$y'=(x^2+y^2)^{1/2}$
\item(1)\label{exer:2.3.10}$y'=x(y^2-1)^{2/3}$
\item(1)\label{exer:2.3.11}$y'=(x^2+y^2)^2$
\item(1)\label{exer:2.3.12}$y'=(x+y)^{1/2}$


\item\label{exer:2.3.13}$y'=\dfrac{\tan y }{x-1}$

\item\label{exer:2.3.14}Apply \zcref{thmtype:2.3.1} to the initial value problem
\[
y'+p(x)y = q(x), \quad y(x_0)=y_0
\]
for a linear equation, and compare the conclusions that can be drawn
from it to those that follow from Theorem 2.1.2.

\item\label{exer:2.3.15}\leavevmode\vspace{-\baselineskip}
\begin{alist}
\item %(a)
Verify that the function
\[
y =
\begin{cases}
(x^2-1)^{5/3}, & -1 < x < 1, \\
0, & |x| \ge 1,
\end{cases}
\]
is a solution of the initial value problem
\[
y'=\frac{10}{3}xy^{2/5}, \quad y(0)=-1
\]
on $(-\infty,\infty)$. \hint{You'll need the definition
\[
y'(\overline{x})
= \lim_{x \to\overline{x}}\frac{y(x)-y(\overline{x})}{x-\overline{x}}
\]
to verify that $y$ satisfies the differential
equation at $\overline{x} = \pm 1$.}

\item %(b)
Verify that if $\epsilon_i=0$ or $1$ for $i=1$, $2$ and $a$, $b>1$, then
the function
\[
y =
\begin{cases}
\epsilon_1(x^2-a^2)^{5/3}, & - \infty < x < -a, \\
0, & -a \le x \le -1, \\
(x^2-1)^{5/3}, &  -1 < x < 1, \\
0, & 1 \le x \le b, \\
\epsilon_2(x^2-b^2)^{5/3}, & b < x < \infty,
\end{cases}
\]
is a solution of the initial value problem of \part{a} on
$(-\infty,\infty)$.
\end{alist}

\item\label{exer:2.3.16}Use the ideas developed in \zcref{exer:2.3.15} to find
infinitely
many solutions  of the initial value problem
\[
y'=y^{2/5}, \quad y(0)=1
\]
on $(-\infty,\infty)$.


\item\label{exer:2.3.17}Consider the initial value
problem
\[
 y' = 3x(y-1)^{1/3}, \quad y(x_0) = y_0.\tag{A}
\]

\begin{alist}
\item %(a)
For what points $(x_0,y_0)$ does \zcref{thmtype:2.3.1} imply that
(A) has a solution?
\item %(b)
 For what points $(x_0,y_0)$ does \zcref{thmtype:2.3.1} imply that
(A) has a  unique solution on some open interval
that contains  $x_0$?
\end{alist}

\item\label{exer:2.3.18}Find nine solutions of the initial value problem
\[
y'=3x(y-1)^{1/3}, \quad y(0)=1
\]
that are all defined on $(-\infty,\infty)$ and differ from each other
for values of  $x$ in every open interval that contains $x_0=0$.

\item\label{exer:2.3.19}From \zcref{thmtype:2.3.1}, the initial
value
problem
 \[
y'=3x(y-1)^{1/3}, \quad y(0)=9
\]
has a unique solution on an open interval that contains  $x_0=0$. Find the
solution and determine the largest open interval on which it's unique.

\item\label{exer:2.3.20}\leavevmode\vspace{-\baselineskip}
\begin{alist}
\item %(a)
 From \zcref{thmtype:2.3.1},
the initial value problem
\[
 y'=3x(y-1)^{1/3}, \quad y(3)=-7\tag{A}
\]
 has a unique solution on some open interval that contains $x_0=3$.
Determine
the largest such open interval, and find the solution on this interval.
\item %(b)
 Find infinitely many solutions of (A), all defined
on
$(-\infty,\infty)$.
\end{alist}


\item\label{exer:2.3.21}Prove:
\begin{alist}
\item % ()
If
\[
f(x,y_0) = 0,\quad a<x<b,\tag{A}
\]
and  $x_0$ is in  $(a,b)$,
then $y\equiv y_0$ is a solution of
\[
y'=f(x,y), \quad y(x_0)=y_0
\]
on $(a,b)$.
\item % (b)
If $f$ and $f_y$ are continuous on an open rectangle
that contains $(x_0,y_0)$ and (A) holds, no
solution of $y'=f(x,y)$ other than $y\equiv y_0$ can equal $y_0$ at
any point in $(a,b)$.
\end{alist}
\end{exerciselist}


\section{Transformation of Nonlinear Equations into Separable Equations}
\label{sec:transform_separable}
\markright{\thesection. \MakeUppercase{Transformation of Nonlinear Equations}}

In \zcref{sec:linar_first_order} we found that the solutions of a linear
nonhomogeneous equation
\[
y'+p(x)y=f(x)
\]
are of the form $y=uy_1$, where $y_1$ is a nontrivial solution of the
complementary equation
\begin{equation} \label{eq:2.4.1}
y'+p(x)y=0
\end{equation}
and $u$ is a solution of
\[
u'y_1(x)=f(x).
\]
Note that this last equation is separable, since it can be rewritten
as
\[
u'=\frac{f(x)}{y_1(x)}.
\]
In this section we'll consider nonlinear differential equations
that are not separable to begin with, but can be solved in a similar
fashion by writing their solutions in the form $y=uy_1$, where $y_1$
is a suitably chosen known function and $u$ satisfies a separable
equation. We'llsay in this case  that we {\color{blue}\it
transformed\/} the given equation into a separable equation.

\boxit{Bernoulli Equations}

A \href{http://www-history.mcs.st-and.ac.uk/Mathematicians/Bernoulli_Jacob.html}
{\color{blue}\itshape Bernoulli
 equation} is an equation  of
the form
\begin{equation} \label{eq:2.4.2}
y'+p(x)y=f(x)y^r,
\end{equation}
where $r$ can be any real number other than $0$ or $1$. (Note that
\eqref{eq:2.4.2} is linear if and only if $r=0$ or $r=1$.) We can
transform
\eqref{eq:2.4.2} into a separable equation by variation of parameters:
if  $y_1$ is  a nontrivial solution of \eqref{eq:2.4.1},
substituting
$y=uy_1$ into \eqref{eq:2.4.2} yields
\[
u'y_1+u(y_1'+p(x)y_1)=f(x)(uy_1)^r,
\]
which is equivalent to the separable equation
\[
u'y_1(x)=f(x)\left(y_1(x)\right)^ru^r\mbox{\quad or \quad}
\frac{u'}{u^r}=f(x)\left(y_1(x)\right)^{r-1},
\]
since $y_1'+p(x)y_1=0$.

\begin{example}\label{example:2.4.1}
Solve the Bernoulli equation
\begin{equation} \label{eq:2.4.3}
y'-y=xy^2.
\end{equation}
\end{example}

\solution
Since $y_1=e^x$ is a solution of $y'-y=0$, we look for solutions of
\eqref{eq:2.4.3}  in the form $y=ue^x$, where
\[
u'e^x=xu^2e^{2x}\mbox{\quad or, equivalently, \quad}
u'=xu^2e^x.
\]
Separating variables and integrating yields
\begin{align*}
\frac{u'}{u^2}&=xe^x\\
-\frac{1}{u}&=(x-1)e^x+c\\
u&=-\frac{1}{(x-1)e^x+c}\\
y&=-\frac{1}{x-1+ce^{-x}}.
\end{align*}
%
\zcref[S]{figure:2.4.1}  shows
direction field and some integral curves of \eqref{eq:2.4.3}.

\begin{figure}
  \centering
  \includegraphics[alt={A vector field for -2<x,y<2 with arrows pointing right and downward for negative x and right and upward for positive x.  Curves below a horizontal asymptote y=0 decrease to bottom out near x=0.5 and then tend towards (2,-0.5).  Curves above y=0 bottom out near x=-1 and then increase toward vertical asymptotes.},
  width=5.67in,height=3.3in,keepaspectratio]{fig020401}
\color{blue}
\caption{A direction field and integral curves for $y'-y=xy^{2}$}
  \label{figure:2.4.1}
\end{figure}


\boxit{Other Nonlinear Equations That Can be  Transformed Into
Separable Equations}

We've seen that the nonlinear Bernoulli equation can be transformed
into a separable equation by  the substitution $y=uy_1$ if
$y_1$ is suitably chosen. Now let's discover a sufficient condition
for a nonlinear first order differential equation
\begin{equation} \label{eq:2.4.4}
y'=f(x,y)
\end{equation}
to be transformable into a separable equation in the same way.
  Substituting $y=uy_1$  into
\eqref{eq:2.4.4} yields
\[
u'y_1(x)+uy_1'(x)=f(x,uy_1(x)),
\]
which is equivalent to
\begin{equation} \label{eq:2.4.5}
u'y_1(x)=f(x,uy_1(x))-uy_1'(x).
\end{equation}
If
\[
f(x,uy_1(x))=q(u)y_1'(x)
\]
for some function $q$, then   \eqref{eq:2.4.5} becomes
\begin{equation} \label{eq:2.4.6}
u'y_1(x)=(q(u)-u)y_1'(x),
\end{equation}
which is separable. After checking for constant solutions $u\equiv
u_0$ such that $q(u_0)=u_0$, we can separate
variables to obtain
\[
\frac{u'}{q(u)-u}=\frac{y_1'(x)}{y_1(x)}.
\]

\boxit{Homogeneous Nonlinear Equations}

In the text  we'll consider only the most widely studied class of
equations for which the method of the preceding paragraph works.
Other types of equations appear in
\zcref[range]{exer:2.4.44,exer:2.4.51}.

The differential equation \eqref{eq:2.4.4}
is said to be \emph{\color{blue}homogeneous} if  $x$ and $y$
occur in $f$ in such a way that  $f(x,y)$ depends only
on the ratio $y/x$; that is, \eqref{eq:2.4.4} can be written as
\begin{equation} \label{eq:2.4.7}
y'=q(y/x),
\end{equation}
where $q=q(u)$ is a function of a single variable.
For example,
\[
y'=\frac{y+xe^{-y/x}}{x}=\frac{y}{x}+e^{-y/x}
\]
and
\[
y'=\frac{y^2+xy-x^2}{x^2}=\left(\frac{y}{x}\right)^2+\frac{y}{x}
-1
\]
are of the form  \eqref{eq:2.4.7}, with
\[
q(u)=u+e^{-u}\quad\text{and}\quad  q(u)=u^2+u-1,
\]
respectively. The general method discussed above can be
applied to
\eqref{eq:2.4.7} with $y_1=x$ (and therefore $y_1'=1)$. Thus,
substituting $y=ux$ in \eqref{eq:2.4.7} yields
\[
u'x+u=q(u),
\]
and separation of variables (after checking for constant
solutions $u\equiv u_0$ such that $q(u_0)=u_0$) yields
\[
\frac{u'}{q(u)-u}=\frac{1}{x}.
\]

Before turning to examples, we point out something that you may've have
already noticed:
 the definition of \emph{\color{blue}homogeneous equation} given
here isn't  the same as the definition given in \zcref{sec:linar_first_order},
where we said that a linear equation of the form
\[
y'+p(x)y=0
\]
is homogeneous. We make no apology for this inconsistency, since we
didn't create it     historically, \emph{\color{blue}homogeneous} has been
used in these two inconsistent ways. The one
having to do with linear equations is the most important. This
is the only section of the book where the meaning defined here will
apply.

Since $y/x$ is in general undefined if $x=0$, we'll consider
solutions of nonhomogeneous equations only on open intervals that do
not contain the point $x=0$.

\begin{example}\label{example:2.4.2}
Solve
\begin{equation} \label{eq:2.4.8}
y'=\frac{y+xe^{-y/x}}{x}.
\end{equation}
\end{example}

\solution Substituting  $y=ux$
into \eqref{eq:2.4.8} yields
\[
u'x+u =  \frac{ux+xe^{-ux/x}}{x} = u+e^{-u}.
\]
 Simplifying and separating variables yields
\[
e^uu'=\frac{1}{x}.
\]
Integrating yields
$e^u=\ln |x|+c$.
Therefore
$u=\ln(\ln|x|+c)$ and
$y=ux=x \ln (\ln |x|+c)$.

\zcref[S]{figure:2.4.2} shows
a direction field and integral curves for \eqref{eq:2.4.8}.

\begin{figure}
  \centering
  \includegraphics[alt={A vector field for 0.5<x<5 and -2<y<2 with arrows pointing up and right.  The arrows are more vertical for small x values, and point more northeast as x increases.},
  width=5.67in,height=3.3in,keepaspectratio]{fig020402}
\color{blue}
\caption
{A direction field and some integral curves for
$y'=\dfrac{y+xe^{-y/x}}{x}$}
  \label{figure:2.4.2}
\end{figure}

\begin{example}\label{example:2.4.3}
\begin{alist}
\item %(a)
Solve
\begin{equation} \label{eq:2.4.9}
x^2y'=y^2+xy-x^2.
\end{equation}

\item %(b)
Solve the initial value problem
\begin{equation} \label{eq:2.4.10}
x^2y'=y^2+xy-x^2, \quad y(1)=2.
\end{equation}
\end{alist}
\end{example}

\begin{alistSolution}
\item
We first find solutions of \eqref{eq:2.4.9} on open intervals that don't
contain $x=0$. We can rewrite \eqref{eq:2.4.9} as
\[
y'=\frac{y^2+xy-x^2}{x^2}
\]
for $x$ in any such interval. Substituting $y=ux$ yields
\[
u'x+u =\frac{ (ux)^2+x(ux)-x^2}{x^2}
= u^2+u-1,
\]
so
\begin{equation} \label{eq:2.4.11}
u'x=u^2-1.
\end{equation}
By inspection this equation has the constant solutions $u\equiv1$ and
$u\equiv-1$. Therefore $y=x$ and $y=-x$ are solutions of
\eqref{eq:2.4.9}. If $u$ is a solution of \eqref{eq:2.4.11} that doesn't
assume the values $\pm 1$ on some interval,  separating variables
yields
\[
\frac{u'}{u^2-1}=\frac{1}{x},
\]
 or, after a partial fraction expansion,
\[
\frac{1}{2}\left[\frac{1}{u-1}-\frac{1}{u+1}\right]u'=
\frac{1}{x}.
\]
 Multiplying by 2 and integrating yields
\[
\ln\left|\frac{u-1}{u+1}\right| =2 \ln |x|+k,
\]
 or
\[
\left|\frac{u-1}{u+1}\right|=e^kx^2,
\]
which holds if
\begin{equation} \label{eq:2.4.12}
\frac{u-1}{u+1}=cx^2
\end{equation}
where $c$ is an arbitrary constant.
  Solving for $u$ yields
\[
u =\frac{1+cx^2}{1-cx^2}.
\]
%
\begin{minipage}[t]{0.48\linewidth}
\tagstructbegin{tag=Sect}
\includegraphics[alt={A vector field for -2<x,y<2 with straight line solutions x=0 going upward, and y=±x going to the right. Elsewhere, the arrows smoothly interpolate the straight line solutions.  Several solution curves follow the arrows.},
width=\linewidth,height=2.44in,keepaspectratio]{fig020403}
\captionof{figure}{A direction field and integral curves for
$x^{2}y'=y^{2}+xy-x^{2}$}
\label{figure:2.4.3}
\tagstructend
\end{minipage}%
\hfill
\begin{minipage}[t]{0.48\linewidth}
\tagstructbegin{tag=Sect}
\includegraphics[alt={Several solution curves from the previous vector field for x<0.  For x>0, the solution curves have all coalesced.},
width=\linewidth,height=2.44in,keepaspectratio]{fig020404}
\captionof{figure}{Solutions of $x^{2}y'=y^{2}+xy-x^{2}$,\; $y(1)=2$}
\label{figure:2.4.4}
\tagstructend
\end{minipage}\bigskip\\
%
Therefore
\begin{equation} \label{eq:2.4.13}
y=ux=\frac{x(1+cx^2)}{1-cx^2}
\end{equation}
is a solution of \eqref{eq:2.4.10} for any choice of the constant $c$.
Setting $c=0$ in \eqref{eq:2.4.13} yields the solution $y=x$. However, the
solution $y=-x$ can't be obtained from \eqref{eq:2.4.13}. Thus, the
solutions of \eqref{eq:2.4.9} on intervals that don't contain $x=0$ are
$y=-x$ and functions of the form \eqref{eq:2.4.13}.

The situation is more complicated if $x=0$ is in the open interval.
 First, note that $y=-x$ satisfies \eqref{eq:2.4.9}
on $(-\infty,\infty)$. If $c_1$ and $c_2$ are arbitrary constants,
 the function
\begin{equation} \label{eq:2.4.14}
y=\begin{cases}
\dfrac{x(1+c_1x^2)}{1-c_1x^2},&a<x<0,\\
\dfrac{x(1+c_2x^2)}{1-c_2x^2},&0\le x<b,
\end{cases}
\end{equation}
is a solution of \eqref{eq:2.4.9} on $(a,b)$, where
\[
a=\begin{cases}
-\dfrac{1}{\sqrt{c_1}}&\text{ if }c_1>0,\\
-\infty&\text{ if }c_1\le 0,
\end{cases}
\quad\text{and}\quad
b=\begin{cases}
\dfrac{1}{\sqrt{c_2}}&\text{ if }c_2>0,\\
\infty&\text{ if }c_2\le 0.
\end{cases}
\]
We leave it to you to verify this. To do so, note that if $y$ is
any function of the form \eqref{eq:2.4.13} then $y(0)=0$ and $y'(0)=1$.


\zcref[S]{figure:2.4.3} shows a direction field and some integral curves
for \eqref{eq:2.4.9}.

\item  We could obtain $c$ by imposing
the initial condition $y(1)=2$ in \eqref{eq:2.4.13}, and then solving for
$c$. However, it's easier to use \eqref{eq:2.4.12}. Since $u=y/x$, the
initial
condition
$y(1)=2$ implies that $u(1)=2$.  Substituting this into \eqref{eq:2.4.12}
yields $c=1/3$.  Hence, the solution of \eqref{eq:2.4.10} is
\[
y=\frac{x(1+x^2/3)}{1-x^2/3}.
\]
The interval of validity of this solution is $(-\sqrt3,\sqrt3)$.
However, the largest interval on which \eqref{eq:2.4.10} has a unique
solution is $(0,\sqrt3)$. To see this, note from \eqref{eq:2.4.14}
that any function of the form
\begin{equation} \label{eq:2.4.15}
y=
\begin{cases}
\dfrac{x(1+cx^2)}{1-cx^2},&a<x\le0,\\
\dfrac{x(1+x^2/3)}{1-x^2/3},&0\le x<\sqrt3,
\end{cases}
\end{equation}
is a solution of \eqref{eq:2.4.10} on $(a,\sqrt3)$, where $a=-1/\sqrt c$
if $c>0$ or $a=-\infty$ if $c\le0$. (Why doesn't this contradict
\zcref{thmtype:2.3.1}?)


\zcref[S]{figure:2.4.4} shows several solutions of the
initial value problem~\eqref{eq:2.4.10}. Note that these solutions coincide
on $(0,\sqrt{3})$.
\end{alistSolution}

In the last two examples we were able to solve the given equations
explicitly.   However, this isn't  always possible, as you'll
see in the exercises.


\exercises
\emph{In \zcref[range]{exer:2.4.1,exer:2.4.4} solve the given Bernoulli
equation.}

\begin{exerciselist}

\item(1)\label{exer:2.4.1}$y'+y=y^2$
\item(1)\label{exer:2.4.2}$7xy'-2y=-\dfrac{x^2 }{y^6}$
\item(1)\label{exer:2.4.3}$x^2y'+2y=2e^{1/x}y^{1/2}$
\item(1)\label{exer:2.4.4}$(1+x^2)y'+2xy=\frac{1}{(1+x^2)y}$

\exercisetext{In \zcref[range]{exer:2.4.5,exer:2.4.6} find all
solutions. Also, plot a direction field and some integral curves on
the indicated rectangular region.}

 \item\label{exer:2.4.5}\CGex
$y'-xy=x^3y^3;    \quad \{-3\le x\le 3,\-2\le y\ge 2\}$

\item\label{exer:2.4.6}\CGex
$y'-\dfrac{1+x}{3x}y=y^4;   \quad \{-2\le x\le2,-2\le y \le2\}$

\exercisetext{In \zcref[range]{exer:2.4.7,exer:2.4.11} solve the
initial value problem.}

\item\label{exer:2.4.7}$y'-2y=xy^3,\quad y(0)=2\sqrt2$

\item\label{exer:2.4.8}$y'-xy=xy^{3/2},\quad y(1)=4$
\item\label{exer:2.4.9}$xy'+y=x^4y^4,\quad y(1)=1/2$
\item\label{exer:2.4.10}$y'-2y=2y^{1/2},\quad y(0)=1$
\item\label{exer:2.4.11}$y'-4y=\dfrac{48x}{y^2},\quad y(0)=1$



\exercisetext{In \zcref[range]{exer:2.4.12,exer:2.4.13} solve the
initial value problem and graph the solution.}

\item\label{exer:2.4.12}\CGex $x^2y'+2xy=y^3,\quad y(1)=1/\sqrt2$
\item\label{exer:2.4.13}\CGex $y'-y=xy^{1/2},\quad y(0)=4$


\item\label{exer:2.4.14}You may have noticed that  the logistic equation
\[
P'=aP(1-\alpha P)
\]
from Verhulst's model for population growth can be written in
Bernoulli form as
\[
P'-aP=-a\alpha P^2.
\]
This isn't particularly interesting, since the logistic equation is
separable, and therefore solvable by the method studied in
\zcref{sec:separable}. So let's consider a more complicated model,
where
$a$ is a positive constant and $\alpha$ is a positive continuous
function of $t$ on $[0,\infty)$. The equation for this model is
\[
P'-aP=-a\alpha(t) P^2,
\]
a non-separable Bernoulli equation.
\begin{alist}
\item % (a)
Assuming that $P(0)=P_0>0$, find $P$ for $t>0$. \hint{Express your
result in terms of the integral
 $\int_0^t\alpha(\tau)e^{a\tau}\,d\tau$.}
\item % (b)
Verify that your result reduces to the known results
for the Malthusian model where  $\alpha=0$, and the Verhulst model
where $\alpha$ is a nonzero constant.
\item % (c)
Assuming that
\[
\lim_{t\to\infty}e^{-at}\int_0^t\alpha(\tau)e^{a\tau}\,d\tau=L
\]
exists (finite or infinite), find $\lim_{t\to\infty}P(t)$.
\end{alist}

\exercisetext{In \zcref[range]{exer:2.4.15,exer:2.4.18} solve
the equation explicitly.}

\item(1)\label{exer:2.4.15}$y'=\dfrac{y+x}{x}$
\item(1)\label{exer:2.4.16}$y'=\dfrac{y^2+2xy }{x^2}$
\item(1)\label{exer:2.4.17}$xy^3y'=y^4+x^4$
\item(1)\label{exer:2.4.18}$y'=\dfrac{y}{x}+\sec\dfrac{y}{x}$


\exercisetext{In \zcref[range]{exer:2.4.19,exer:2.4.21} solve the
equation explicitly.
 Also, plot a direction field and some integral curves on
the indicated rectangular region.}

  \item\label{exer:2.4.19}\CGex
$x^2y'=xy+x^2+y^2;   \quad \{-8\le x\le 8,-8\le y\le 8\}$

 \item\label{exer:2.4.20}\CGex
$xyy'=x^2+2y^2;  \quad \{-4\le x\le 4,-4\le y\le 4\}$


\item\label{exer:2.4.21}\CGex
$y'=\dfrac{2y^2+x^2e^{-(y/x)^2}}{2xy};\quad  \{-8\le x\le 8,-8\le y\le 8\}$


\exercisetext{In \zcref[range]{exer:2.4.22,exer:2.4.27} solve
the initial value problem.}

\item\label{exer:2.4.22}$y'=\dfrac{xy+y^2}{x^2},\quad y(-1)=2$

 \item\label{exer:2.4.23}$y'=\dfrac{x^3+y^3}{xy^2}, \quad y(1)=3$

 \item\label{exer:2.4.24}$xyy'+x^2+y^2=0, \quad y(1)=2$

\item\label{exer:2.4.25}$y'=\dfrac{y^2-3xy-5x^2 }{x^2}, \quad y(1)=-1$

\item\label{exer:2.4.26}$x^2y'=2x^2+y^2+4xy, \quad y(1)=1$

\item\label{exer:2.4.27}$xyy'=3x^2+4y^2, \quad y(1)=\sqrt{3}$

\exercisetext{In \zcref[range]{exer:2.4.28,exer:2.4.34} solve the
given homogeneous equation implicitly.}

\item(1)\label{exer:2.4.28}$y'=\dfrac{x+y }{x-y}$
\item(1)\label{exer:2.4.29}$(y'x-y)(\ln |y|-\ln |x|)=x$
\item(1)\label{exer:2.4.30}$y'=\dfrac{y^3+2xy^2+x^2y+x^3}{x(y+x)^2}$ 
\item(1)\label{exer:2.4.31}$y'=\dfrac{x+2y }{2x+y}$
\item(1)\label{exer:2.4.32}$y'=\dfrac{y }{y-2x}$
\item(1)\label{exer:2.4.33}$y'=\dfrac{xy^2+2y^3}{x^3+x^2y+xy^2}$

\item\label{exer:2.4.34}$y'=\dfrac{x^3+x^2y+3y^3 }{x^3+3xy^2}$


\item\label{exer:2.4.35}\Lex
\begin{alist}
\item % (a)
Find a solution of the initial value problem
\[
x^2y'=y^2+xy-4x^2, \quad y(-1)=0\tag{A}
\]
on the interval $(-\infty,0)$. Verify that this solution is actually
valid on $(-\infty,\infty)$.
\item % (b)
Use \zcref{thmtype:2.3.1} to show that (A) has a
unique solution on $(-\infty,0)$.
\item % (c)
Plot a direction field for the differential equation in
(A) on a square
\[
\{-r\le x\le r, -r\le y\le r\},
\]
where $r$
is any positive number. Graph the solution you obtained in
\part{a} on this field.
\item % (d)
Graph  other solutions  of (A) that are defined on $(-\infty,\infty)$.
\item % (e)
Graph other solutions of (A)  that are defined only on intervals of
the form $(-\infty,a)$, where is a finite positive number.
\end{alist}



\item\label{exer:2.4.36}\Lex
\begin{alist}
\item % (a)
Solve the equation
\[
xyy'=x^2-xy+y^2\tag{A}
\]
implicitly.
\item % (b)
Plot a direction field for (A) on a square
\[
\{0\le x\le r,0\le y\le r\}
\]
where $r$ is any positive number.
\item % (c)
Let $K$ be a positive integer. (You may have to try several
choices for $K$.)
Graph solutions of the initial value problems
\[
xyy'=x^2-xy+y^2,\quad y(r/2)=\frac{kr}{K},
\]
for $k=1$, $2$, \dots, $K$. Based on your observations, find
 conditions on the positive numbers $x_0$
and $y_0$ such that the initial value problem
\[
xyy'=x^2-xy+y^2,\quad y(x_0)=y_0,\tag{B}
\]
has  a unique solution (i) on $(0,\infty)$   or (ii) only
on an interval  $(a,\infty)$, where $a>0$?
\item % (d)
What can you say about the graph of the solution of (B)
as $x\to\infty$? (Again, assume that $x_0>0$ and $y_0>0$.)
\end{alist}


\item\label{exer:2.4.37}\Lex
\begin{alist}
\item % (a)
Solve the equation
\[
y'=\frac{2y^2-xy+2x^2 }{xy+2x^2}\tag{A}
\]
implicitly.
\item % (b)
Plot a direction field for (A) on a square
\[
\{-r\le x\le r,-r\le y\le r\}
\]
where $r$ is any positive number. By graphing solutions of (A),
determine necessary and sufficient conditions on $(x_0,y_0)$ such that
(A) has a solution on (i) $(-\infty,0)$ or (ii) $(0,\infty)$ such that
$y(x_0)=y_0$.
\end{alist}

\item\label{exer:2.4.38}\Lex
Follow the instructions of \zcref{exer:2.4.37} for the equation
\[
y'=\frac{xy+x^2+y^2 }{xy}.
\]

\item\label{exer:2.4.39}\Lex
Pick any  nonlinear homogeneous equation $y'=q(y/x)$  you like, and
plot  direction fields on the square $\{-r\le x\le r,\ -r\le y\le
r\}$,
where $r>0$. What happens to the direction field as you vary $r$?
Why?

\item\label{exer:2.4.40}Prove:  If $ad-bc\ne 0$, the equation
\[
y'=\frac{ax+by+\alpha }{cx+dy+\beta}
\]
can be transformed into the homogeneous nonlinear equation
\[
\frac{dY }{dX}=\frac{aX+bY }{cX+dY}
\]
by the substitution $x=X-X_0,\  y=Y-Y_0$,
where $X_0$ and $Y_0$ are suitably chosen constants.

\exercisetext{In \zcref[range]{exer:2.4.41,exer:2.4.43} use a
method suggested by \zcref{exer:2.4.40} to solve the
given equation implicitly.}

\item(1)\label{exer:2.4.41}$y'=\dfrac{-6x+y-3 }{2x-y-1}$
\item(1)\label{exer:2.4.42}$y'=\dfrac{2x+y+1 }{x+2y-4}$

\item\label{exer:2.4.43}$y'=\dfrac{-x+3y-14 }{x+y-2}$

\exercisetext{In \zcref[range]{exer:2.4.44,exer:2.4.51} find
a function $y_1$ such that the substitution $y=uy_1$ transforms
the given equation  into a separable
equation of the form \eqref{eq:2.4.6}. Then solve the given equation
explicitly.}

\item(1)\label{exer:2.4.44}$3xy^2y'=y^3+x$
\item(1)\label{exer:2.4.45}$xyy'=3x^6+6y^2$
\item(1)\label{exer:2.4.46}$x^3y'=2(y^2+x^2y-x^4)$
\item(1)\label{exer:2.4.47}$y'=y^2e^{-x}+4y+2e^x$
\item(1)\label{exer:2.4.48}$y'=\dfrac{y^2+y\tan x+\tan^2 x}{\sin^2x}$ 
\item(1)\label{exer:2.4.49}$x(\ln x)^2y'=-4(\ln x)^2+y\ln x+y^2$
\item(1)\label{exer:2.4.50}$2x(y+2\sqrt x)y'=(y+\sqrt x)^2$
\item(1)\label{exer:2.4.51}$(y+e^{x^2})y'=2x(y^2+ye^{x^2}+e^{2x^2})$

\medskip
\item\label{exer:2.4.52}Solve the initial value problem
\[
y'+\frac{2}{x}y=\frac{3x^2y^2+6xy+2}{x^2(2xy+3)},\quad y(2)=2.
\]

\item\label{exer:2.4.53}Solve the initial value problem
\[
y'+\frac{3}{x}y=\frac{3x^4y^2+10x^2y+6}{x^3(2x^2y+5)},\quad y(1)=1.
\]

\item\label{exer:2.4.54}Prove:  If $y$ is a solution of a homogeneous nonlinear equation
$y'=q(y/x)$,  so is $y_1=y(ax)/a$, where $a$ is any nonzero
constant.


\item\label{exer:2.4.55}A \emph{\color{blue}generalized}
\href{http://http://www-history.mcs.st-and.ac.uk/Indexes/Riccati.html}
{\color{blue}\itshape Riccati equation} is of the form
\[
y'=P(x)+Q(x)y+R(x)y^2.\tag{A}
\]
(If $R\equiv-1$,  (A) is a
\href{http://http://www-history.mcs.st-and.ac.uk/Indexes/Riccati.html}
{\color{blue}\itshape Riccati
equation\/}.) Let $y_1$ be a known solution and $y$ an arbitrary
solution of (A). Let $z=y-y_1$. Show that $z$ is a
solution of a Bernoulli equation with $n=2$.


\exercisetext{In \zcref[range]{exer:2.4.56,exer:2.4.59},
given that $y_1$ is a solution of the given equation, use the
method suggested by \zcref{exer:2.4.55} to find other solutions.}

\item\label{exer:2.4.56}$y'=1+x - (1+2x)y+xy^2$;  \quad    $y_1=1$

\item\label{exer:2.4.57}$y'=e^{2x}+(1-2e^x)y+y^2$;  \quad    $y_1=e^x$

\item\label{exer:2.4.58}$xy'=2-x+(2x-2)y-xy^2$;  \quad    $y_1=1$

\item\label{exer:2.4.59}$xy'=x^3+(1-2x^2)y+xy^2$;  \quad    $y_1=x$

\end{exerciselist}


\section{Exact Equations}\label{sec:exact}

In this  section it's convenient to write first order
differential equations in the form
\begin{equation} \label{eq:2.5.1}
M(x,y)\,dx+N(x,y)\,dy=0.
\end{equation}
This equation  can be interpreted as
\begin{equation} \label{eq:2.5.2}
M(x,y)+N(x,y)\,\frac{dy}{dx}=0,
\end{equation}
where $x$ is the independent variable and $y$ is the dependent
variable, or as
\begin{equation} \label{eq:2.5.3}
M(x,y)\,\frac{dx}{dy}+N(x,y)=0,
\end{equation}
where $y$ is the independent variable and $x$ is the dependent
variable. Since the solutions of \eqref{eq:2.5.2} and \eqref{eq:2.5.3} will
often have to be left in
implicit, form we'll say that $F(x,y)=c$ is an implicit solution of
\eqref{eq:2.5.1} if every differentiable function $y=y(x)$ that satisfies
$F(x,y)=c$ is a solution of \eqref{eq:2.5.2} and every
differentiable function $x=x(y)$ that satisfies $F(x,y)=c$ is a
solution of \eqref{eq:2.5.3}.

Here are  some examples:

\begin{flalign*}\toprule
\text{\zcref{eq:2.5.1}}&&
\text{\zcref{eq:2.5.2}}&&
\text{\zcref{eq:2.5.3}}
\\\midrule
3x^2y^2\,dx+2x^3y\,dy &=0 & 3x^2y^2+2x^3y\,\dfrac{dy}{dx} &=0 & 3x^2y^2\,\dfrac{dx}{dy}+2x^3y &=0 \\
(x^2+y^2)\,dx +2xy\,dy&=0 &
(x^2+y^2)+2xy\,\dfrac{dy}{dx}&=0&
(x^2+y^2)\,\dfrac{dx}{dy} +2xy&=0\\
3y\sin x\,dx-2xy\cos x\,dy &=0 &
3y\sin x-2xy\cos x\,\dfrac{dy}{dx} &=0 &
3y\sin x\,\dfrac{dx}{dy}-2xy\cos x  &=0\\\bottomrule
\end{flalign*}\stepcounter{table}

Note that a separable equation can be written as
\eqref{eq:2.5.1} as
\[
M(x)\,dx+N(y)\,dy=0.
\]


We'll  develop a method for solving \eqref{eq:2.5.1} under appropriate
assumptions on $M$ and $N$. This method is an extension
of the method of separation of variables
(\zcref{exer:2.5.41}).  Before stating it we
consider an  example.

\begin{example}\label{example:2.5.1}
Show that
\begin{equation} \label{eq:2.5.4}
x^4y^3+x^2y^5+2xy=c
\end{equation}
is an implicit solution of
\begin{equation} \label{eq:2.5.5}
(4x^3y^3+2xy^5+2y)\,dx+(3x^4y^2+5x^2y^4+2x)\,dy=0.
\end{equation}
\end{example}

\solution
Regarding $y$ as a function of $x$ and
differentiating \eqref{eq:2.5.4} implicitly with respect to
$x$ yields
\[
(4x^3y^3+2xy^5+2y)+(3x^4y^2+5x^2y^4+2x)\,\frac{dy}{dx}=0.
\]
Similarly, regarding $x$ as a function of $y$ and
differentiating \eqref{eq:2.5.4} implicitly with respect to
$y$ yields
\[
(4x^3y^3+2xy^5+2y)\frac{dx}{dy}+(3x^4y^2+5x^2y^4+2x)=0.
\]
Therefore \eqref{eq:2.5.4} is an implicit solution of \eqref{eq:2.5.5}
in either of its two possible interpretations. \bbox

You may think this example is pointless, since
concocting a differential equation that has a given implicit solution
isn't particularly interesting. However, it illustrates the
next important theorem, which  we'll prove by using implicit
differentiation,  as  in  \zcref{example:2.5.1}.

\begin{theorem}\label{thmtype:2.5.1}
If $F=F(x,y)$ has continuous partial derivatives
$F_x$ and $F_y$, then
\begin{equation} \label{eq:2.5.6}
F(x,y)=c
\end{equation}
(where $c$ is a constant) is an implicit solution of the differential equation
\begin{equation} \label{eq:2.5.7}
F_x(x,y)\,dx+F_y(x,y)\,dy=0.
\end{equation}
\end{theorem}

\begin{proof}
Regarding $y$ as a function of $x$ and  differentiating
\eqref{eq:2.5.6}  implicitly with respect to $x$ yields
\[
F_x(x,y)+F_y(x,y)\,\frac{dy}{dx}=0.
\]
On the other hand,
 regarding $x$ as a function of $y$ and  differentiating
\eqref{eq:2.5.6}  implicitly with respect to $y$ yields
\[
F_x(x,y)\,\frac{dx}{dy}+F_y(x,y)=0.
\]
Thus, \eqref{eq:2.5.6} is an
implicit solution of  \eqref{eq:2.5.7} in either of its two possible
interpretations.
\end{proof}

We'll say that  the equation
\begin{equation} \label{eq:2.5.8}
M(x,y)\,dx+N(x,y)\,dy=0
\end{equation}
 is  \emph{\color{blue}exact} on an an open rectangle  $R$ if there's
a function $F=F(x,y)$ such  $F_x$
and $F_y$  are continuous, and
\begin{equation} \label{eq:2.5.9}
F_x(x,y)=M(x,y) \quad\text{and}\quad  F_y(x,y)=N(x,y)
\end{equation}
for  all  $(x,y)$ in $R$.
This usage of ``exact'' is related  to its usage in calculus,
where the expression
\[
F_x(x,y)\,dx+F_y(x,y)\,dy
\]
(obtained by substituting \eqref{eq:2.5.9} into the left side of
\eqref{eq:2.5.8}) is the  \emph{\color{blue}exact differential of} $F$.

\zcref[S]{example:2.5.1} shows that it's easy to solve
\eqref{eq:2.5.8}
if it's exact \emph{\color{blue}and}  we know a function $F$ that satisfies
\eqref{eq:2.5.9}. The important questions are:
\begin{enumext}[
 label=\arabic*,
 font=\scshape,
 wrap-label={Question #1.},
 labelwidth=6em,
 ]
 \item Given an equation \eqref{eq:2.5.8}, how can we determine whether it's  exact?
\item If \eqref{eq:2.5.8} is exact, how do we find
a function $F$ satisfying \eqref{eq:2.5.9}?
\end{enumext}

To discover the answer to Question~1,
 assume that  there's a function $F$  that satisfies \eqref{eq:2.5.9} on
some open rectangle $R$, and in addition that $F$ has continuous mixed
partial derivatives $F_{xy}$ and $F_{yx}$.  Then a theorem from calculus
implies that
\begin{equation} \label{eq:2.5.10}
F_{xy}=F_{yx}.
\end{equation}
If $F_x=M$ and $F_y=N$,
differentiating the first of these equations with respect to
$y$ and the second with respect to $x$ yields
\begin{equation} \label{eq:2.5.11}
F_{xy}=M_y\quad\text{and}\quad   F_{yx}=N_x.
\end{equation}
From  \eqref{eq:2.5.10}  and \eqref{eq:2.5.11}, we conclude that
 a necessary condition for exactness is that $M_y=N_x$.
This motivates the next theorem, which we state without proof.

\begin{theorem}[The Exactness Condition]\label{thmtype:2.5.2}
Suppose  $M$ and
$N$ are continuous and have continuous partial derivatives
$M_y$ and $N_x$ on an open rectangle $R.$ Then
\[
M(x,y)\,dx+N(x,y)\,dy=0
\]
is exact on $R$ if and only if
\begin{equation} \label{eq:2.5.12}
M_y(x,y)=N_x(x,y)
\end{equation}
for all $(x,y)$ in  $R.$.
\end{theorem}

To help you  remember the exactness condition, observe
that the coefficients of $dx$ and $dy$ are differentiated in
\eqref{eq:2.5.12} with respect to the ``opposite'' variables; that is,
the coefficient of $dx$ is differentiated with respect to $y$, while the
coefficient of $dy$ is differentiated with respect to $x$.

\begin{example}\label{example:2.5.2}
 Show that the equation
\[
3x^2y\,dx+4x^3\,dy=0
\]
is not exact on any open rectangle.
\end{example}

\solution   Here
\[
M(x,y)=3x^2y\quad\text{and}\quad  N(x,y)=4x^3
\]
so
\[
M_y(x,y)=3x^2 \quad\text{and}\quad  N_x(x,y)=12 x^2.
\]
Therefore  $M_y=N_x$   on the line $x=0$,
but not on any open rectangle, so
there's no
function $F$ such that $F_x(x,y)=M(x,y)$ and $F_y(x,y)=N(x,y)$
for all $(x,y)$ on any open rectangle.
\bbox

The next example illustrates two possible methods for finding a
function $F$ that satisfies the condition $F_x=M$ and $F_y=N$ if
$M\,dx+N\,dy=0 $ is exact.

\begin{example}\label{example:2.5.3}
Solve
\begin{equation}\label{eq:2.5.13}
(4x^3y^3+3x^2)\,dx+(3x^4y^2+6y^2)\,dy=0.
\end{equation}
\end{example}


\solution
Here
\[
M(x,y)=4x^3y^3+3x^2,\quad N(x,y)=3x^4y^2+6y^2,
\]
and
\[
M_y(x,y)=N_x(x,y)=12 x^3y^2
\]
for all $(x,y)$.
Therefore
\zcref{thmtype:2.5.2} implies that there's a function $F$ such that
\begin{equation} \label{eq:2.5.14}
F_x(x,y)=M(x,y)=4x^3y^3+3x^2
\end{equation}
 and
\begin{equation} \label{eq:2.5.15}
F_y(x,y)=N(x,y)=3x^4y^2+6y^2
\end{equation}
 for all $(x,y)$.

(Method 1)  To find $F$, we integrate \eqref{eq:2.5.14} with
respect to $x$ to obtain
\begin{equation} \label{eq:2.5.16}
F(x,y)=x^4y^3+x^3+\phi(y),
\end{equation}
 where $\phi (y)$ is the ``constant'' of integration.  (Here
$\phi$ is ``constant'' in  that it's independent of $x$, the
variable of integration.)  If $\phi$ is any differentiable function of
$y$ then $F$  satisfies \eqref{eq:2.5.14}.  To
determine $\phi$ so that
$F$ also satisfies \eqref{eq:2.5.15}, assume that $\phi$ is
differentiable and differentiate $F$ with respect to $y$.
This yields
\[
F_y(x,y)=3x^4y^2+\phi'(y).
\]
 Comparing this with \eqref{eq:2.5.15} shows that
\[
\phi'(y)=6y^2.
\]
 We integrate this with respect to $y$ and take the
constant of integration to be zero because we're interested only in
finding \emph{\color{blue}some} $F$ that satisfies \eqref{eq:2.5.14} and
\eqref{eq:2.5.15}. This  yields
\[
\phi (y)=2y^3.
\]
Substituting this into \eqref{eq:2.5.16} yields
\begin{equation} \label{eq:2.5.17}
F(x,y)=x^4y^3+x^3+2y^3.
\end{equation}
Now \zcref{thmtype:2.5.1} implies that
\[
x^4y^3+x^3+2y^3=c
\]
is an implicit solution of \eqref{eq:2.5.13}. Solving this for $y$
yields the explicit solution
\[
y=\left(\frac{c-x^3}{2+x^4}\right)^{1/3}.
\]

(Method 2) Instead of first integrating
\eqref{eq:2.5.14}
with respect to $x$, we could begin by integrating \eqref{eq:2.5.15} with
respect to $y$ to obtain
\begin{equation} \label{eq:2.5.18}
F(x,y)=x^4y^3+2y^3+\psi (x),
\end{equation}
 where $\psi$ is an arbitrary  function of
$x$.  To determine $\psi$, we assume that $\psi$ is
differentiable and differentiate $F$ with respect to $x$,
which yields
\[
F_x(x,y)=4x^3y^3+\psi'(x).
\]
 Comparing this with \eqref{eq:2.5.14} shows that
\[
\psi'(x)=3x^2.
\]
Integrating this and again taking  the constant of
integration to be zero yields
\[
\psi(x)=x^3.
\]
 Substituting this into \eqref{eq:2.5.18} yields \eqref{eq:2.5.17}.

\zcref[S]{figure:2.5.1} shows a direction field and  some
integral curves of \eqref{eq:2.5.13}.

\begin{figure}
  \centering
  \includegraphics[alt={A vector field for -1<x,y<1 with arrows flowing down and to the right. Several solution curves follow the arrows, coalescing as they approach ±(1,0.7).  Along y=x, the curves bump away from the origin.},
  width=5.67in,height=3.3in,keepaspectratio]{fig020501}
\color{blue}
\caption{A direction field and  integral curves for
$(4x^3y^3+3x^2)\,dx+(3x^4y^2+6y^2)\,dy=0$}
  \label{figure:2.5.1}
\end{figure}



Here's a summary of the procedure used in Method 1 of this
example. You should summarize procedure used in Method~2.

\smallskip
{\color{blue}
{\centerline{\bfseries Procedure For Solving An Exact Equation}}

\begin{enumext}[
label=\arabic*,
wrap-label={Step #1.},
labelwidth={4em},
list-indent=0pt,
]
\item Check that the equation
\begin{equation} \label{eq:2.5.19}
 M(x,y)\,dx+N(x,y)\,dy=0
\end{equation}
satisfies the exactness condition  $M_y=N_x$. If not,
don't go further with this procedure.
\item Integrate
\[
\frac{\partial F(x,y)}{\partial x}=M(x,y)
\]
with respect to $x$ to obtain
\begin{equation} \label{eq:2.5.20}
F(x,y)=G(x,y)+\phi(y),
\end{equation}
where $G$ is an antiderivative of $M$ with respect to $x$, and $\phi$
is an unknown function of $y$.
\item Differentiate \eqref{eq:2.5.20} with respect to
$y$  to obtain
\[
\frac{\partial F(x,y)}{\partial y}=\frac{\partial G(x,y)}{\partial
y}+\phi'(y).
\]
\item Equate the right side of this equation to $N$ and solve
for $\phi'$;    thus,
\[
\frac{\partial G(x,y)}{\partial y}+\phi'(y)=N(x,y),\mbox{\quad  so \quad }
\phi'(y)=N(x,y)-\frac{\partial G(x,y)}{\partial y}.
\]
\item Integrate $\phi'$ with respect to $y$, taking the
constant of integration to be zero, and substitute the result in
\eqref{eq:2.5.20} to obtain $F(x,y)$.
\item Set $F(x,y)=c$ to obtain an implicit solution of
\eqref{eq:2.5.19}. If possible, solve for $y$ explicitly as a
function of $x$.
\end{enumext}
}


It's a common mistake to omit Step 6. However, it's important to
include this step, since $F$ isn't  itself a solution of
\eqref{eq:2.5.19}.

Many equations can be conveniently solved by either of the two methods
used in \zcref{example:2.5.3}. However, sometimes the integration
required in one approach is  more difficult than in the
other. In such cases we choose the approach that requires the easier
integration.

\begin{example}\label{example:2.5.4}
 Solve the equation
\begin{equation}  \label{eq:2.5.21}
(ye^{xy} \tan x+e^{xy} \sec^2x)\,dx+xe^{xy} \tan x\,dy=0.
\end{equation}
\end{example}

\solution We leave it to you to check that
$M_y=N_x$ on any open rectangle  where $\tan x$ and $\sec x$ are
defined. Here we must find a function $F$ such that
\begin{equation} \label{eq:2.5.22}
F_x(x,y)=ye^{xy} \tan x+e^{xy} \sec^2 x
\end{equation}
 and
\begin{equation}  \label{eq:2.5.23}
F_y(x,y)=xe^{xy} \tan x.
\end{equation}
 It's difficult to integrate \eqref{eq:2.5.22} with respect to $x$, but
easy to integrate \eqref{eq:2.5.23} with respect to $y$.  This yields
\begin{equation} \label{eq:2.5.24}
F(x,y)=e^{xy} \tan x+\psi (x).
\end{equation}
Differentiating this with respect to $x$ yields
\[
F_x(x,y)=ye^{xy}\tan x+e^{xy}\sec^2x+\psi'(x).
\]
 Comparing this with \eqref{eq:2.5.22} shows that  $\psi'(x)=0$.
Hence, $\psi$ is  a constant, which we can take to be zero in
\eqref{eq:2.5.24}, and
\[
e^{xy} \tan x=c
\]
is an implicit solution of \eqref{eq:2.5.21}. \bbox

Attempting to apply our procedure to an equation that isn't  exact
will lead to failure in  Step 4, since the function
\[
 N-\frac{\partial G}{\partial y}
\]
won't be independent of $x$ if $M_y\ne N_x$
(\zcref{exer:2.5.31}), and therefore
can't be the derivative of a function of $y$ alone. Here's an example that
illustrates this.

\begin{example}\label{example:2.5.5}
 Verify that the equation
\begin{equation} \label{eq:2.5.25}
3x^2y^2\,dx+6x^3y\,dy=0
\end{equation}
is not exact, and show  that the procedure for solving exact equations
  fails when applied to \eqref{eq:2.5.25}.
\end{example}

\solution   Here
\[
M_y(x,y)=6x^2y \mbox{\quad  and \quad } N_x(x,y)=18x^2y,
\]
so \eqref{eq:2.5.25} isn't  exact. Nevertheless,  let's try to find
a function $F$ such that
\begin{equation} \label{eq:2.5.26}
F_x(x,y)=3x^2y^2
\end{equation}
and
\begin{equation} \label{eq:2.5.27}
F_y(x,y)=6x^3y.
\end{equation}
Integrating \eqref{eq:2.5.26} with respect to $x$ yields
\[
F(x,y)=x^3y^2+\phi(y),
\]
and differentiating this with respect to $y$ yields
\[
F_y(x,y)=2x^3y+\phi'(y).
\]
For this equation to be consistent with \eqref{eq:2.5.27},
\[
6x^3y=2x^3y+\phi'(y),
\]
or
\[
\phi'(y)=4x^3y.
\]
This is a contradiction, since $\phi'$  must be independent
of $x$. Therefore the procedure fails.

\exercises
\emph{In \zcref[range]{exer:2.5.1,exer:2.5.17} determine which equations
are exact and solve them.}

\begin{exerciselist}
\item\label{exer:2.5.1}$6x^2y^2\,dx+4x^3y\,dy=0$

\item\label{exer:2.5.2}$(3y\cos x+4xe^x+2x^2e^x)\,dx+(3\sin x+3)\,dy=0$

\item\label{exer:2.5.3}$14x^2y^3\,dx+21 x^2y^2\,dy=0$

\item\label{exer:2.5.4}$(2x-2y^2)\,dx+(12y^2-4xy)\,dy=0$

% these were two columns
\item\label{exer:2.5.5}$(x+y)^2\,dx+(x+y)^2\,dy=0$%&
\item\label{exer:2.5.6}$(4x+7y)\,dx+(3x+4y)\,dy=0$

\item\label{exer:2.5.7}$(-2y^2\sin x+3y^3-2x)\,dx+(4y\cos x+9xy^2)\,dy=0$

\item\label{exer:2.5.8}$(2x+y)\,dx+(2y+2x)\,dy=0$

\item\label{exer:2.5.9}$(3x^2+2xy+4y^2)\,dx+(x^2+8xy+18y)\,dy=0$

\item\label{exer:2.5.10}$(2x^2+8xy+y^2)\,dx+(2x^2+xy^3/3)\,dy=0$

\item\label{exer:2.5.11}$\left(\dfrac{1}{x}+2x\right)\,dx+
\left(\dfrac{1}{y}+2y\right)\,dy=0$

\item\label{exer:2.5.12}$(y\sin xy+xy^2\cos xy)\,dx+(x\sin xy+xy^2\cos xy)\,dy=0$

\item\label{exer:2.5.13}$\dfrac{x\,dx}{(x^2+y^2)^{3/2}}+\dfrac{y\,dy}{(x^2+y^2)^{3/2}}=0$

\item\label{exer:2.5.14}$\left(e^x(x^2y^2+2xy^2)+6x\right)\,dx+(2x^2ye^x+2)\,dy=0$

\item\label{exer:2.5.15}$\left(x^2e^{x^2+y}(2x^2+3)+4x\right)\,dx+(x^3e^{x^2+y}-12y^2)\,dy=0$

\item\label{exer:2.5.16}$\left(e^{xy}(x^4y+4x^3)+3y\right)\,dx+(x^5e^{xy}+3x)\,dy=0$

\item\label{exer:2.5.17}$(3x^2\cos xy-x^3y\sin xy+4x)\,dx+(8y-x^4\sin xy)\,dy=0$

\exercisetext{In \zcref[range]{exer:2.5.18,exer:2.5.22} solve
the initial value problem.}

\item\label{exer:2.5.18}$(4x^3y^2-6x^2y-2x-3)\,dx+(2x^4y-2x^3)\,dy=0,\quad y(1)=3$

\item\label{exer:2.5.19}$(-4y\cos x+4\sin x\cos x+\sec^2x)\,dx+
(4y-4\sin x)\,dy=0,\quad y(\pi/4)=0$

\item\label{exer:2.5.20}$(y^3-1)e^x\,dx+3y^2(e^x+1)\,dy=0,\quad y(0)=0$

\item\label{exer:2.5.21}$(\sin x-y\sin x-2\cos x)\,dx+\cos x\,dy=0,\quad y(0)=1$

\item\label{exer:2.5.22}$(2x-1)(y-1)\,dx+(x+2)(x-3)\,dy=0,\quad y(1)=-1$

\item\label{exer:2.5.23}\CGex
Solve the exact equation
\[
(7x+4y)\,dx+(4x+3y)\,dy=0.
\]
Plot a direction field and some integral curves for this equation on
the rectangle
\[
\{-1\le x\le1,-1\le y\le1\}.
\]

\item\label{exer:2.5.24}\CGex
Solve the exact equation
\[
e^x(x^4y^2+4x^3y^2+1)\,dx+(2x^4ye^x+2y)\,dy=0.
\]
Plot a direction field and some integral curves for this equation on
the rectangle
\[
\{-2\le x\le2,-1\le y\le1\}.
\]

\item\label{exer:2.5.25}\CGex
Plot a direction field and some integral curves for the exact equation
\[
(x^3y^4+x)\,dx+(x^4y^3+y)\,dy=0
\]
on the rectangle $\{-1\le x\le 1,-1\le y\le1\}$. (See
\zcref{exer:2.5.37}(a)).

\item\label{exer:2.5.26}\CGex
Plot a direction field and some integral curves for the exact equation
\[
(3x^2+2y)\,dx+(2y+2x)\,dy=0
\]
on the rectangle $\{-2\le x\le 2,-2\le y\le2\}$. (See
\zcref{exer:2.5.37}(b)).

\item\label{exer:2.5.27}\Lex
\begin{alist}
\item % (a)
Solve the exact equation
\[
(x^3y^4+2x)\,dx+(x^4y^3+3y)\,dy=0\tag{A}
\]
implicitly.
\item % (b)
For what choices of $(x_0,y_0)$ does
\zcref{thmtype:2.3.1} imply that the initial value problem
\[
(x^3y^4+2x)\,dx+(x^4y^3+3y)\,dy=0,\quad y(x_0)=y_0,\tag{B}
\]
has a unique solution on an open interval $(a,b)$  that contains $x_0$?
\item % (b)
Plot a direction field and some integral curves for (A)
on a rectangular region centered at the origin. What is the
 interval of validity of the solution of (B)?
\end{alist}

\item\label{exer:2.5.28}\Lex
\begin{alist}
\item % (a)
Solve the exact equation
\[
(x^2+y^2)\,dx+2xy\,dy=0\tag{A}
\]
implicitly.
\item % (b)
For what choices of $(x_0,y_0)$ does
\zcref{thmtype:2.3.1} imply that the initial value problem
\[
(x^2+y^2)\,dx+2xy\,dy=0,\quad y(x_0)=y_0,\tag{B}
\]
has a unique solution $y=y(x)$ on some open interval $(a,b)$
that contains $x_0$?
\item % (c)
Plot a direction field and some integral curves for (A). From
the plot determine,
the  interval $(a,b)$ of \part{b}, the monotonicity
properties (if any) of the solution of (B), and $\lim_{x\to
a+}y(x)$ and $\lim_{x\to b-}y(x)$. \hint{Your answers will
depend upon which quadrant contains $(x_0,y_0)$.}
\end{alist}

\item\label{exer:2.5.29}Find all functions $M$ such that the  equation is exact.
\begin{alist}
\item %(a)
$M(x,y)\,dx+(x^2-y^2)\,dy=0$
\item %(b)
$M(x,y)\,dx+2xy\sin x\cos y\,dy=0$
\item %(c)
$M(x,y)\,dx+(e^x-e^y\sin x)\,dy=0$
\end{alist}

\item\label{exer:2.5.30}Find all functions $N$ such that the  equation is exact.
\begin{alist}
\item %(a)
$(x^3y^2+2xy+3y^2)\,dx+N(x,y)\,dy=0$
\item %(b)
$(\ln xy+2y\sin x)\,dx+N(x,y)\,dy=0$
\item %(c)
$(x\sin x+y\sin y)\,dx+N(x,y)\,dy=0$
\end{alist}

\item\label{exer:2.5.31}Suppose $M,N,$ and their partial derivatives are continuous on
an open rectangle
$R$, and $G$ is an antiderivative of $M$ with respect to $x$; that is,
\[
\frac{\partial G}{\partial x}=M.
\]
Show that if $M_y\ne N_x$ in $R$ then the function
\[
 N-\frac{\partial  G}{\partial y}
\]
is not independent of $x$.

\item\label{exer:2.5.32}Prove:  If the equations $M_1\,dx+N_1\,dy=0$ and $M_2\,
dx+N_2\,dy=0$ are exact on an open rectangle $R$,  so is
the equation \[(M_1+M_2)\,dx+(N_1+N_2)\,dy=0.\]

\item\label{exer:2.5.33}Find conditions on the constants $A$, $B$, $C$, and $D$ such that
the equation
\[
(Ax+By)\,dx+(Cx+Dy)\,dy=0
\]
is exact.

\item\label{exer:2.5.34}Find conditions on the constants $A$, $B$, $C$, $D$, $E$, and
$F$ such that the equation
\[
(Ax^2+Bxy+Cy^2)\,dx+(Dx^2+Exy+Fy^2)\,dy=0
\]
is exact.

\item\label{exer:2.5.35}Suppose $M$ and $N$ are continuous and have continuous partial
derivatives $M_y$ and $N_x$ that satisfy the exactness condition
$M_y=N_x$ on an open rectangle $R$.
  Show that if $(x,y)$ is in $R$ and
\[
F(x,y)=\int^x_{x_0}M(s,y_0)\,ds+\int^y_{y_0}N(x,t)\,dt,
\]
then $F_x=M$ and $F_y=N$.

\item\label{exer:2.5.36}Under the assumptions of \zcref{exer:2.5.35}, show that
\[
F(x,y)=\int^y_{y_0}N(x_0,s)\,ds+\int^x_{x_0}M(t,y)\,dt.
\]

\item\label{exer:2.5.37}Use the method suggested by \zcref{exer:2.5.35}, with
$(x_0,y_0)=(0,0)$, to solve the these exact equations:
\begin{alist}
\item %(a)
$(x^3y^4+x)\,dx+(x^4y^3+y)\,dy=0$
\item %(b)
$(x^2+y^2)\,dx+2xy\,dy=0$
\item %(c)
$(3x^2+2y)\,dx+(2y+2x)\,dy=0$
\end{alist}

\item\label{exer:2.5.38}Solve the initial value problem
\[
y'+\frac{2}{x}y=-\frac{2xy}{x^2+2x^2y+1},\quad y(1)=-2.
\]

\item\label{exer:2.5.39}Solve the initial value problem
\[
y'-\frac{3}{x}y=\frac{2x^4(4x^3-3y)}{3x^5+3x^3+2y},\quad y(1)=1.
\]

\item\label{exer:2.5.40}Solve the initial value problem
\[
y'+2xy=-e^{-x^2}\left(\frac{3x+2ye^{x^2}}{2x+3ye^{x^2}}\right),\quad
y(0)=-1.
\]

\item\label{exer:2.5.41}Rewrite the separable equation
\[
h(y)y'=g(x)\tag{A}
\]
as an exact equation
\[
M(x,y)\,dx+N(x,y)\,dy=0.\tag{B}
\]
Show that applying the method of this section  to
(B) yields the same solutions that would be
obtained by applying the method of separation of variables
to~(A)

\item\label{exer:2.5.42}Suppose  all second partial derivatives of $M=M(x,y)$ and $N=N(x,y)$
are continuous and $M\,dx+N\,dy=0$ and $-N\,dx+M\,dy=0$ are
exact on an open rectangle  $R$.  Show that
$M_{xx}+M_{yy}=N_{xx}+N_{yy}=0$ on
$R$.

\item\label{exer:2.5.43}Suppose all second partial derivatives of
$F=F(x,y)$ are continuous and $F_{xx}+F_{yy}=0$ on an open rectangle
$R$. (A function with these properties is said to be \emph{\color{blue}harmonic};
see also \zcref{exer:2.5.42}.) Show that $-F_y\,dx+F_x\,dy=0$ is
exact on $R$, and
therefore there's a function $G$ such that
$G_x=-F_y$ and $G_y=F_x$ in $R$. (A function $G$ with this property is
said to be a \emph{\color{blue}harmonic conjugate} of $F$.)

\item\label{exer:2.5.44}Verify that the following functions are harmonic, and find all their
harmonic conjugates.  (See \zcref{exer:2.5.43}.)

\begin{alist}[columns=3]
% TODO horizontal
\item $x^2-y^2$
\item $e^x\cos y$
\item $x^3-3xy^2$
\item $\cos x\cosh y$
\item $\sin x\cosh y$
\end{alist}

\end{exerciselist}


\section{Integrating Factors}\label{sec:integrating_factors}

In \zcref{sec:exact} we saw that if $M$, $N$, $M_y$ and $N_x$ are
continuous and $M_y=N_x$ on an open rectangle $R$ then
\begin{equation} \label{eq:2.6.1}
M(x,y)\,dx+N(x,y)\,dy=0
\end{equation}
is exact on $R$. Sometimes an equation that isn't  exact can be made
exact by multiplying it by an appropriate function. For example,
\begin{equation}\label{eq:2.6.2}
(3x+2y^2)\,dx+2xy\,dy=0
\end{equation}
is  not exact, since
$M_y(x,y)=4y\ne  N_x(x,y)=2y$ in \eqref{eq:2.6.2}.
 However, multiplying \eqref{eq:2.6.2}  by $x$ yields
\begin{equation}\label{eq:2.6.3}
(3x^2+2xy^2)\,dx+2x^2y\,dy=0,
\end{equation}
which is exact, since
$M_y(x,y)=N_x(x,y)=4xy$ in \eqref{eq:2.6.3}.
Solving \eqref{eq:2.6.3} by the procedure given in \zcref{sec:exact}
yields the implicit solution
\[
x^3+x^2y^2=c.
\]

A function $\mu=\mu(x,y)$ is  an \emph{\color{blue}integrating factor} for
\eqref{eq:2.6.1}  if
\begin{equation}\label{eq:2.6.4}
 \mu(x,y)M (x,y)\,dx+\mu(x,y)N (x,y)\,dy=0
 \end{equation}
 is exact. If we know an integrating
factor $\mu$ for \eqref{eq:2.6.1}, we can solve the exact equation
\eqref{eq:2.6.4} by the method of \zcref{sec:exact}. It would be
nice
if we could say that \eqref{eq:2.6.1} and \eqref{eq:2.6.4} always have the
same solutions, but this isn't so. For example, a solution
$y=y(x)$ of \eqref{eq:2.6.4} such that $\mu(x,y(x))=0$ on some interval
$a<x<b$ could fail to be a solution of \eqref{eq:2.6.1}
(\zcref{exer:2.6.1}), while
\eqref{eq:2.6.1} may have a solution $y=y(x)$ such that $\mu(x,y(x)) $
isn't even defined (\zcref{exer:2.6.2}). Similar comments
apply if $y$ is the independent variable and $x$ is the dependent
variable  in \eqref{eq:2.6.1} and \eqref{eq:2.6.4}.  However, if $\mu(x,y)$
is defined and nonzero for all $(x,y)$,  \eqref{eq:2.6.1}  and
\eqref{eq:2.6.4} are equivalent; that is, they have the same solutions.

\boxit{Finding Integrating Factors}

By applying \zcref{thmtype:2.5.2} (with $M$ and $N$ replaced by $\mu M$
and $\mu N$), we see that \eqref{eq:2.6.4} is exact on an open rectangle
$R$ if $\mu M$, $\mu N$, $(\mu M)_y$, and $(\mu N)_x$ are continuous
and
\[
\frac{\partial}{\partial y}(\mu M)=\frac{\partial}{\partial x}
(\mu N) \mbox{\quad or, equivalently, \quad}
\mu_yM+\mu M_y=\mu_xN+\mu N_x
\]
on $R$. It's better to rewrite the last equation as
\begin{equation} \label{eq:2.6.5}
\mu(M_y-N_x)=\mu_xN-\mu_yM,
\end{equation}
which  reduces to  the known result for exact equations;
that is, if $M_y=N_x$ then \eqref{eq:2.6.5} holds with $\mu=1$, so
\eqref{eq:2.6.1} is exact.

You may think  \eqref{eq:2.6.5} is of little value, since it involves
\emph{\color{blue}partial} derivatives of the unknown integrating factor $\mu$,
and we haven't studied methods for solving such equations. However,
we'll now show that \eqref{eq:2.6.5} is useful if we restrict our search
to integrating factors that are products of a function of
$x$ and a function of $y$; that is,
$\mu(x,y)=P(x)Q(y)$.
We're not saying that \emph{\color{blue}every} equation
 $M\,dx+N\,dy=0$ has an integrating factor of this
form;   rather, we're saying that \emph{\color{blue}some} equations have
such integrating
factors.We'llnow develop a way to determine whether a given
equation has such an integrating factor, and a method for finding the
integrating factor in this case.


If $\mu(x,y)=P(x)Q(y)$, then $\mu_x(x,y)=P'(x)Q(y)$ and
$\mu_y(x,y)=P(x)Q'(y)$, so \eqref{eq:2.6.5} becomes
\begin{equation} \label{eq:2.6.6}
P(x)Q(y)(M_y-N_x)=P'(x)Q(y)N-P(x)Q'(y)M,
\end{equation}
or, after dividing through by $P(x)Q(y)$,
\begin{equation} \label{eq:2.6.7}
M_y-N_x=\frac{P'(x)}{P(x)}N-\frac{Q'(y)}{Q(y)}M.
\end{equation}
Now let
\[
p(x)=\frac{P'(x)}{P(x)} \quad\text{and}\quad  q(y)=\frac{Q'(y)}{Q(y)},
\]
so \eqref{eq:2.6.7} becomes
\begin{equation} \label{eq:2.6.8}
M_y-N_x=p(x)N-q(y)M.
\end{equation}

We obtained \eqref{eq:2.6.8} by \emph{\color{blue}assuming} that $M\,dx+N\,dy=0$ has
an integrating factor $\mu(x,y)=P(x)Q(y)$. However, we can
now view \eqref{eq:2.6.7} differently: If there are functions $p=p(x)$
and $q=q(y)$  that satisfy \eqref{eq:2.6.8} and we define
\begin{equation} \label{eq:2.6.9}
P(x)=\pm e^{\int p(x)\,dx}\quad\text{and}\quad 
Q(y)=\pm e^{\int q(y)\,dy},
\end{equation}
then reversing the steps that led from \eqref{eq:2.6.6} to
\eqref{eq:2.6.8} shows that $\mu(x,y)=P(x)Q(y)$ is an integrating factor
for $M\,dx+N\,dy=0$. In using this result, we take the constants of
integration in \eqref{eq:2.6.9} to be zero and choose the signs
conveniently so the integrating factor has the simplest form.

There's no simple general method for ascertaining
whether functions $p=p(x)$ and $q=q(y)$ satisfying \eqref{eq:2.6.8}
exist. However, the next theorem gives simple sufficient
conditions for the given equation to have an integrating factor that
depends on only one of the independent variables $x$ and $y$, and for
finding an integrating factor in this case.

\begin{theorem}\label{thmtype:2.6.1}
Let $M,$ $N,$ $M_y,$ and $N_x$ be continuous on an open rectangle $R.$
Then$:$
\begin{alist}
\item % (a)
 If $(M_y-N_x)/N$ is independent
of $y$ on $R$ and  we define
\[
p(x)=\frac{M_y-N_x}{N}
\]
 then
\begin{equation}\label{eq:2.6.10}
\mu(x)=\pm e^{\int p(x)\,dx}
\end{equation}
is an integrating factor for
\begin{equation}\label{eq:2.6.11}
M(x,y)\,dx+N(x,y)\,dy=0
\end{equation}
on $R.$
\item % (b)
If $(N_x-M_y)/M$ is independent of $x$ on
$R$ and we define
\[
q(y)=\frac{N_x-M_y}{M},
\]
then
\begin{equation}\label{eq:2.6.12}
\mu(y)=\pm e^{\int q(y)\,dy}
\end{equation}
 is an integrating factor for \eqref{eq:2.6.11} on $R.$
\end{alist}
\end{theorem}

\medskip

\begin{proof}
(a)
If $(M_y-N_x)/N$ is independent of $y$, then \eqref{eq:2.6.8}
 holds with $p=(M_y-N_x)/N$ and $q\equiv0$. Therefore
 \[
P(x)=\pm e^{\int p(x)\,dx}\mbox{\quad and\quad}Q(y)=\pm e^{\int
q(y)\,dy}=\pm e^0=\pm1,
 \]
so \eqref{eq:2.6.10} is an integrating factor for \eqref{eq:2.6.11} on $R$.

(b) If $(N_x-M_y)/M$ is independent of $x$ then eqref{eq:2.6.8} holds
with $p\equiv0$ and $q=(N_x-M_y)/M$, and a similar argument shows that
\eqref{eq:2.6.12} is an integrating factor for \eqref{eq:2.6.11} on $R$.
\end{proof}

The next two examples show how to apply
\zcref{thmtype:2.6.1}.

\begin{example}\label{example:2.6.1}
Find an integrating factor for the equation
\begin{equation}\label{eq:2.6.13}
(2xy^3-2x^3y^3-4xy^2+2x)\,dx+(3x^2y^2+4y)\,dy=0
\end{equation}
and solve the equation.
\end{example}

\solution In \eqref{eq:2.6.13}
\[
M=2xy^3-2x^3y^3-4xy^2+2x,\ N=3x^2y^2+4y,
\]
and
\[
M_y-N_x=(6xy^2-6x^3y^2-8xy)-6xy^2=-6x^3y^2-8xy,
\]
so \eqref{eq:2.6.13}  isn't  exact. However,
\[
\frac{M_y-N_x}{N}=-\frac{6x^3y^2+8xy}{3x^2y^2+4y}=-2x
\]
is independent of $y$, so \zcref{thmtype:2.6.1}\part{a} applies
with
$p(x)=-2x$. Since
\[
\int p (x)\,dx=-\int 2x\,dx=-x^2,
\]
 $\mu(x)=e^{-x^2}$ is an
integrating factor.  Multiplying \eqref{eq:2.6.13} by $\mu$ yields the
exact equation
\begin{equation}\label{eq:2.6.14}
e^{-x^2}(2xy^3-2x^3y^3-4xy^2+2x)\,dx+
 e^{-x^2}(3x^2y^2+4y)\,dy=0.
\end{equation}

To solve this equation, we must find a
function $F$ such that
\begin{equation}\label{eq:2.6.15}
F_x(x,y)=e^{-x^2}(2xy^3-2x^3y^3-4xy^2+2x)
\end{equation}
 and
\begin{equation}\label{eq:2.6.16}
F_y(x,y)=e^{-x^2}(3x^2y^2+4y).
\end{equation}
 Integrating \eqref{eq:2.6.16} with respect to $y$ yields
\begin{equation}\label{eq:2.6.17}
F(x,y)=e^{-x^2}(x^2y^3+2y^2)+\psi(x).
\end{equation}
 Differentiating this with respect to $x$ yields
\[
F_x(x,y)=e^{-x^2}(2xy^3-2x^3y^3-4xy^2)+\psi'(x).
\]
Comparing this with \eqref{eq:2.6.15} shows that $\psi'(x)=
2xe^{-x^2}$;  therefore, we can let $\psi(x)=-e^{-x^2}$ in
\eqref{eq:2.6.17} and conclude that
\[
e^{-x^2}\left(y^2(x^2y+2)-1\right)=c
\]
is an implicit solution of \eqref{eq:2.6.14}. It is also an implicit solution
of \eqref{eq:2.6.13}.

\zcref[S]{figure:2.6.1}  shows a  direction field and some integral curves
for \eqref{eq:2.6.13}

\begin{figure}
  \centering
  \includegraphics[alt={A vector field for -2<x<2 and -4<y<4 that is symmetric along x=0.  Near the origin, solution curves circle clockwise.  For y>0.5, solution curves decrease and are horizontal near x=0.  For x=-2 and y<-1.3, solution curves increase until x=-1.5 and then decrease toward an asymptote along y=1/x.  For x=-2 and -1.3<y<1, solution curves increase slowly and then around x=-1.5 they turn upward and then track close to y=1.  For y=-4 and 0<x<2, solution curves move upward and then turn left to cross x=0 horizontally.},
  %%% TODO.  WE HAVEN'T SAID HOW A DIRECTION FIELD CAN GO LEFT
  width=5.67in,height=3.3in,keepaspectratio]{fig020601}
\color{blue}
\caption{A direction field and  integral curves for
$(2xy^3-2x^3y^3-4xy^2+2x)\,dx+(3x^2y^2+4y)\,dy=0$}
  \label{figure:2.6.1}
\end{figure}



\begin{example}\label{example:2.6.2}
Find an integrating factor for
\begin{equation}\label{eq:2.6.18}
2xy^3\,dx+(3x^2y^2+x^2y^3+1)\,dy=0
\end{equation}
and solve the equation.
\end{example}

\solution In \eqref{eq:2.6.18},
\[
M=2xy^3,\quad N=3x^2y^2+x^2y^3+1,
\]
and
\[
M_y-N_x=6xy^2-(6xy^2+2xy^3)=-2xy^3,
\]
so \eqref{eq:2.6.18} isn't  exact.  Moreover,
\[
\frac{M_y-N_x}{N}=-\frac{2xy^3}{3x^2y^2+x^2y^2+1}
\]
is not independent of $y$, so \zcref[S]{thmtype:2.6.1}(a) does not
apply. However,   \zcref[S]{thmtype:2.6.1}(b) does apply, since
\[
\frac{N_x-M_y}{M}=\frac{2xy^3}{2xy^3}=1
\]
is independent of $x$, so we can take $q(y)=1$.
 Since
\[
\int q(y)\,dy=\int\,dy=y,
\]
  $\mu(y)=e^y$ is
an integrating factor. Multiplying \eqref{eq:2.6.18} by $\mu$ yields the
exact equation
\begin{equation}\label{eq:2.6.19}
2xy^3e^y\,dx+(3x^2y^2+x^2y^3+1)e^y\,dy=0.
\end{equation}
To solve this equation, we must find a
function $F$ such that
\begin{equation}\label{eq:2.6.20}
F_x(x,y)=2xy^3e^y
\end{equation}
 and
\begin{equation}\label{eq:2.6.21}
F_y(x,y)=(3x^2y^2+x^2y^3+1)e^y.
\end{equation}
 Integrating \eqref{eq:2.6.20} with respect to $x$ yields
\begin{equation}\label{eq:2.6.22}
F(x,y)=x^2y^3e^y+\phi(y).
\end{equation}
Differentiating this with respect to $y$ yields
\[
F_y=(3x^2y^2+x^2y^3)e^y+\phi'(y),
\]
and comparing this with \eqref{eq:2.6.21} shows that $\phi'(y)=e^y$.
 Therefore we set $\phi(y)=e^y$ in \eqref{eq:2.6.22} and conclude
that
\[
(x^2y^3+1)e^y=c
\]
is an implicit solution of \eqref{eq:2.6.19}.
It is also an implicit solution
 of \eqref{eq:2.6.18}. \zcref[S]{figure:2.6.2} shows a direction
field and some integral curves for \eqref{eq:2.6.18}.\bbox

\begin{figure}
  \centering
  \includegraphics[alt={A vector field for -4<x,y<4 symmetric about x=0 with trivial solution y=0. For y>0, solution curves increase from a horizontal asymptote along y=0 and peak horizontally when x=0.  For negative y close to 0, solution curves start near the horizontal asymptote and decrease to a minimum when x=0.  For negative y further from 0, solution curves turn away from x=0 and move downward to the left.},
  width=5.67in,height=3.3in,keepaspectratio]{fig020602}
\color{blue}
\caption{A direction field and integral curves for
$2xy^3e^y\,dx+(3x^2y^2+x^2y^3+1)e^y\,dy=0$}
  \label{figure:2.6.2}
\end{figure}


\zcref[S]{thmtype:2.6.1} does not apply in the next example,
but the more general argument that led to \zcref{thmtype:2.6.1}
provides an integrating factor.

\begin{example}\label{example:2.6.3}
Find an integrating factor for
\begin{equation}\label{eq:2.6.23}
(3xy+6y^2)\,dx+(2x^2+9xy)\,dy=0
\end{equation} and solve the equation.
 \end{example}

\solution In \eqref{eq:2.6.23}
 \[
 M=3xy+6y^2,\ N=2x^2+9xy,
 \]
and
 \[
M_y-N_x=(3x+12y)-(4x+9y)=-x+3y.
\]
 Therefore
 \[ \frac{M_y-N_x}{M}=\frac{-x+3y}{3xy+6y^2} \quad\text{and}\quad  \frac{N_x-M_y}{N}=\frac{x-3y}{2x^2+9xy},
 \]
so \zcref{thmtype:2.6.1} does not apply.
Following the more general argument that led to
\zcref{thmtype:2.6.1}, we look for functions $p=p(x)$ and $q=q(y)$
such that
\[
M_y-N_x=p(x)N-q(y)M;
 \]
 that is,
\[
-x+3y=p(x)(2x^2+9xy)-q(y)(3xy+6y^2).
\]
Since the left side contains
only first degree terms in $x$ and $y$, we rewrite this equation as
\[
xp(x)(2x+9y)-yq(y)(3x+6y)=-x+3y.
 \]
 This will be an identity if
\begin{equation}\label{eq:2.6.24}
 xp(x)=A\mbox{\quad and\quad} yq(y)=B,
\end{equation}
 where $A$ and $B$ are constants such that
\[
-x+3y=A(2x+9y)-B(3x+6y),
 \]
 or, equivalently,
\[
-x+3y=(2A-3B)x+(9A-6B)y.
\]
Equating the coefficients of $x$ and $y$
on both sides shows that the last equation holds for all $(x,y)$ if
\begin{align*}
2A-3B &=-1 \\
9A-6B &=\phantom{-}3,
\end{align*}
 which has the solution $A=1$, $B=1$. Therefore
\eqref{eq:2.6.24} implies that
 \[
 p(x)=\frac{1}{x}\quad\text{and}\quad
q(y)=\frac{1}{y}.
\]
Since
\[
\int p(x)\,dx=\ln|x|\quad\text{and}\quad\int q(y)\,dy=\ln|y|,
\]
 we can let $P(x)=x$ and $Q(y)=y$;
hence, $\mu(x,y)=xy$ is an integrating factor. Multiplying
\eqref{eq:2.6.23} by $\mu$ yields the exact equation
\[
(3x^2y^2+6xy^3)\,dx+(2x^3y+9x^2y^2)\,dy=0.
 \]
 We leave it to you to
use the method of \zcref{sec:exact} to show that this equation has the
implicit solution
\begin{equation}
 \label{eq:2.6.25} x^3y^2+3x^2y^3=c.
\end{equation}
This is also an implicit solution of \eqref{eq:2.6.23}.
Since $x\equiv0$ and $y\equiv0$ satisfy \eqref{eq:2.6.25}, you should
check to see that $x\equiv0$ and $y\equiv0$ are also solutions of
\eqref{eq:2.6.23}. (Why is it necessary to check this?)

\zcref[S]{figure:2.6.3} shows a direction field and integral curves for
\eqref{eq:2.6.23}.


See \zcref{exer:2.6.28} for a general discussion of equations like
\eqref{eq:2.6.23}.


\begin{figure}
  \centering
  \includegraphics[alt={A vector field for -4<x<4 and -2<y<2 with straight line solutions x,y=0 and slant asymptote y=-x/3.  Solutions approach the origin along those lines in the fourth quadrant and alternating with each half line.},
  width=5.67in,height=3.3in,keepaspectratio]{fig020603}
\color{blue}
\caption{A direction field and integral curves for
$(3xy+6y^2)\,dx+(2x^2+9xy)\,dy=0$}
  \label{figure:2.6.3}
\end{figure}


\begin{example}\label{example:2.6.4}
The separable  equation
\begin{equation}\label{eq:2.6.26}
-y\,dx+(x+x^6)\,dy=0
\end{equation}
can be converted to the exact equation
\begin{equation} \label{eq:2.6.27}
-\frac{dx}{x+x^6}+\frac{dy}{y}=0
\end{equation}
by multiplying through
by the  integrating factor
\[
\mu(x,y)=\frac{1}{y(x+x^6)}.
\]
However, to solve \eqref{eq:2.6.27} by the method of \zcref{sec:exact}
we would have to evaluate the nasty integral
\[
\int \frac{dx}{x+x^6}.
\]
Instead, we solve \eqref{eq:2.6.26} explicitly for $y$ by finding  an
integrating factor of the form
$\mu(x,y)=x^ay^b$.
\end{example}

\begin{figure}
  \centering
  \includegraphics[alt={A vector field for -1<x,y<1 symmetric about y=0 with arrows radiating from the origin.  For x near -1 and y>0, the vectors turn upward.  For x near 1 and y>0, the vectors turn more horizontal.},
  width=5.67in,height=3.3in,keepaspectratio]{fig020604}
\color{blue}
\caption{A direction field and integral curves for $-y\,dx+(x+x^6)\,dy=0$}
  \label{figure:2.6.4}
\end{figure}

\solution
In \eqref{eq:2.6.26}
\[
M=-y,\ N=x+x^6,
\]
and
\[
M_y-N_x=-1-(1+6x^5)=-2-6x^5.
\]
We  look for functions
$p=p(x)$ and $q=q(y)$ such that
\[
M_y-N_x=p(x)N-q(y)M;
\]
that is,
\begin{equation}\label{eq:2.6.28}
-2-6x^5=p(x)(x+x^6)+q(y)y.
\end{equation}
The right side will contain the term $-6x^5$ if $p(x)=-6/x$.   Then
\eqref{eq:2.6.28} becomes
\[
-2-6x^5=-6-6x^5+q(y)y,
\]
so  $q(y)=4/y$.  Since
\begin{align*}
\int p(x)\,dx&=-\int\frac{6}{x}\,dx=-6\ln|x|=\ln\frac{1}{x^6},
\qquad\text{and}\\
\int q(y)\,dy&=\phantom-\int\frac{4}{y}\,dy=\phantom-4\ln|y|=\ln{y^4},
\end{align*}
we can take $P(x)=x^{-6}$ and $Q(y)=y^4$,
which yields the integrating factor $\mu(x,y)=x^{-6}y^4$.
Multiplying \eqref{eq:2.6.26} by  $\mu$ yields the exact equation
\[
-\frac{y^5}{x^6}\,dx+\left(\frac{y^4}{x^5}+y^4\right)
\,dy=0.
\]
 We leave it to you to use the method of the \zcref{sec:exact}
to show that this equation has the implicit solution \[
\left(\frac{y}{x}\right)^5+y^5=k.
\]
 Solving for $y$ yields
\[
y=k^{1/5}x(1+x^5)^{-1/5},
\]
which we rewrite as
\[
y=cx(1+x^5)^{-1/5}
\]
by renaming the arbitrary constant.
 This is also a solution of \eqref{eq:2.6.26}.

\zcref[S]{figure:2.6.4} shows a direction field and some integral curves
 for \eqref{eq:2.6.26}.


 \exercises
\begin{exerciselist}

\item\label{exer:2.6.1}\leavevmode\vspace{-\baselineskip}
\begin{alist}
\item % (a)
Verify that $\mu(x,y)=y$ is an integrating factor for
\[
y\,dx+\left(2x+\frac{1}{y}\right)\,dy=0\tag{A}
\]
on any open rectangle that does not intersect  the $x$ axis
or, equivalently,  that
\[
y^2\,dx+(2xy+1)\,dy=0\tag{B}
\]
is exact on any such rectangle.
\item % (b)
Verify that $y\equiv0$  is a solution of  (B),
but not of  (A).
\item % (b)
Show that
\[
y(xy+1)=c\tag{C}
\]
is an implicit solution of (B), and explain why
every differentiable function $y=y(x)$ other than $y\equiv0$
that satisfies  (C)
 is also a solution of~(A).
\end{alist}

\item\label{exer:2.6.2}\leavevmode\vspace{-\baselineskip}
\begin{alist}
\item % (a)
Verify that $\mu(x,y)=1/(x-y)^2$ is an integrating factor for
\[
-y^2\,dx+x^2\,dy=0\tag{A}
\]
on any open rectangle  that does not intersect the line $y=x$
or, equivalently, that
\[
-\frac{y^2}{(x-y)^2}\,dx+\frac{x^2}{(x-y)^2}\,dy=0\tag{B}
\]
is exact on any such rectangle.
\item % (b)
Use \zcref[S]{thmtype:2.2.1} to show that
\[
\frac{xy}{(x-y)}=c\tag{C}
\]
is an implicit solution of (B), and explain why
it's also an implicit solution of (A)
\item % (b)
Verify that $y=x$  is a solution of  (A),
even though it can't be obtained from  (C).
\end{alist}

\exercisetext{In \zcref[range]{exer:2.6.3,exer:2.6.16} find an
integrating factor;
that is a function of only one variable, and solve the given equation.}

\item(1)\label{exer:2.6.3}$y\,dx-x\,dy=0$ 
\item(1)\label{exer:2.6.4}\vspace*{7.5pt} $3x^2y\,dx+2x^3\,dy=0$
\item(1)\label{exer:2.6.5}$2y^3\,dx+3y^2\,dy=0$
\item(1)\label{exer:2.6.6}$(5xy+2y+5)\,dx+2x\,dy=0$

\item\label{exer:2.6.7}$(xy+x+2y+1)\,dx+(x+1)\,dy=0$

\item\label{exer:2.6.8}$(27xy^2+8y^3)\,dx+(18x^2y+12xy^2)\,dy=0$

\item\label{exer:2.6.9}$(6xy^2+2y)\,dx+(12x^2y+6x+3)\,dy=0$

\item\label{exer:2.6.10}$y^2\,dx+\left(xy^2+3xy+\dfrac{1}{y}\right)\,dy=0$


\item\label{exer:2.6.11}$(12x^3y+24x^2y^2)\,dx+(9x^4+32x^3y+4y)\,dy=0$


\item\label{exer:2.6.12}$(x^2y+4xy+2y)\,dx+(x^2+x)\,dy=0$

\item\label{exer:2.6.13}$-y\,dx+(x^4-x)\,dy=0$

\item\label{exer:2.6.14}$\cos x\cos y\,dx
+(\sin x\cos y-\sin x\sin y+y)\,dy=0$

\item\label{exer:2.6.15}$(2xy+y^2)\,dx+(2xy+x^2-2x^2y^2-2xy^3)\,dy=0$

\item\label{exer:2.6.16}$y\sin y\,dx+x(\sin y-y\cos y)\,dy=0$

\exercisetext{In \zcref[range]{exer:2.6.17,exer:2.6.23} find an
integrating
factor of the form $\mu(x,y)=P(x)Q(y)$ and solve the given equation.}


\item\label{exer:2.6.17}$y(1+5\ln|x|)\,dx+4x\ln|x|\,dy=0$

 \item\label{exer:2.6.18}$(\alpha y+ \gamma xy)\,dx+(\beta x+ \delta xy)\,dy=0$

\item\label{exer:2.6.19}$(3x^2y^3-y^2+y)\,dx+(-xy+2x)\,dy=0$

\item\label{exer:2.6.20}$2y\,dx+ 3(x^2+x^2y^3)\,dy=0$

\item\label{exer:2.6.21}$(a\cos xy-y\sin xy)\,dx+(b\cos xy-x\sin xy)\,
dy=0$

\item\label{exer:2.6.22}$x^4y^4\,dx+x^5y^3\,dy=0$

\item\label{exer:2.6.23}$y(x\cos x+2\sin x)\,dx+x(y+1)\sin x\,dy=0$


\exercisetext{In \zcref[range]{exer:2.6.24,exer:2.6.27} find an
integrating factor
and solve the  equation. Plot a direction field and some
integral curves for the equation in the indicated rectangular region.}

\item\label{exer:2.6.24}\CGex
 $(x^4y^3+y)\,dx+(x^5y^2-x)\,dy=0;   \quad \{-1\le x\le1,-1\le y\le1\}$


\item\label{exer:2.6.25}\CGex
$(3xy+2y^2+y)\,dx+(x^2+2xy+x+2y)\,dy=0;   \quad \{-2\le x\le2,-2\le
y\le2\}$


\item\label{exer:2.6.26}\CGex
$(12 xy+6y^3)\,dx+(9x^2+10xy^2)\,dy=0;   \quad \{-2\le x\le2,-2\le y\le2\}$

\item\label{exer:2.6.27}\CGex
$(3x^2y^2+2y)\,dx+ 2x\,dy=0;   \quad \{-4\le x\le4,-4\le y\le4\}$




\item\label{exer:2.6.28}Suppose $a$, $b$, $c$, and $d$ are constants such that
$ad-bc\ne0$, and let $m$ and $n$ be arbitrary real numbers.
Show that
\[
(ax^my+by^{n+1})\,dx+(cx^{m+1}+dxy^n)\,dy=0
\]
has an integrating factor $\mu(x,y)=x^\alpha y^\beta$.

\item\label{exer:2.6.29}Suppose $M$, $N$, $M_x$, and $N_y$  are continuous for all
$(x,y)$, and $\mu=\mu(x,y)$  is an integrating factor for
\[
M(x,y)\,dx+N(x,y)\,dy=0.\tag{A}
\]
Assume that $\mu_x$ and $\mu_y$ are continuous for all $(x,y)$, and
suppose   $y=y(x)$ is a differentiable function such that
$\mu(x,y(x))=0$ and $\mu_x(x,y(x))\ne0$ for all $x$ in some interval
$I$. Show that $y$ is a solution of (A) on $I$.

\item\label{exer:2.6.30}According to \zcref{thmtype:2.1.2}, the general solution of the
linear nonhomogeneous equation
\[
y'+p(x)y=f(x)\tag{A}
\]
is
\[
y=y_1(x)\left(c+\int f(x)/y_1(x)\,dx\right),\tag{B}
\]
where $y_1$ is any nontrivial solution of the
complementary equation $y'+p(x)y=0$. In this exercise we obtain this
conclusion in a different way.
You may find it instructive to apply the method suggested here
to solve some of the exercises in \zcref{sec:linar_first_order}.
\begin{alist}
\item % (a)
Rewrite (A) as
\[
[p(x)y-f(x)]\,dx+\,dy=0,\tag{C}
\]
and show that $\mu=\pm e^{\int p(x)\,dx}$ is an integrating factor
for (C).
\item % (b)
Multiply (A) through by $\mu=\pm e^{\int p(x)\,dx}$  and verify that
the resulting equation can be rewritten as
\[
(\mu(x)y)'=\mu(x)f(x).
\]
Then integrate both sides of this equation and solve for $y$
to show that the general solution of (A) is
\[
y=\frac{1}{\mu(x)}\left(c+\int f(x)\mu(x)\,dx\right).
\]
Why is this form of the general solution
 equivalent to (B)?
\end{alist}

\end{exerciselist}
